\documentclass[11pt,a4paper]{article}
\usepackage[margin=23mm]{geometry}
\usepackage[T1]{fontenc}
\usepackage{lmodern,microtype,amsmath,amssymb,amsthm,mathtools,enumitem,xcolor,booktabs}
\usepackage[colorlinks=true,linkcolor=blue!45!black,citecolor=blue!45!black,urlcolor=blue!45!black]{hyperref}
\newtheorem{theorem}{Theorem}[section]
\newtheorem{lemma}[theorem]{Lemma}
\newtheorem{proposition}[theorem]{Proposition}
\newtheorem{corollary}[theorem]{Corollary}
\theoremstyle{remark}\newtheorem{remark}[theorem]{Remark}
\setlength{\parskip}{3pt}
\setlist{nosep}
\allowdisplaybreaks[1]
\title{Thresholds, coherent structures, and dynamics\\of critical focusing dispersive systems\\[5pt]
\large Kinetic amplification, nonlinear departure,\\and conserved internal correlations}
\author{Shuxing Yang}
\date{8 October 2026}
\begin{document}
\maketitle
\begin{abstract}
We connect sharp critical interpolation with actual nonlinear evolution.
For the pure mass-critical biharmonic equation, zero-momentum
negative-energy solutions at mass excess
$d=(M/M_c)^{4/N}-1>0$ satisfy
$\sup_{t<T_+}\|\Delta u(t)\|_2^2\ge c|E|d^{-5/3}$ for sufficiently
small $d$, in a conservative strong flow with the continuation property.
The proof imposes no radial symmetry or spatial moment: quantitative
constant-phase control produces a slowly moving core, energy splitting
bounds exterior mass, and averaging disjoint annuli retains the favorable
exterior kinetic term in the virial identity. Amplitude perturbations of
sharp standing waves must consequently make arbitrarily large kinetic
excursions as their mass approaches the threshold.
For slow traveling waves we prove constrained energy descent without
real-profile nondegeneracy. In dimensions one through three, with that
nondegeneracy explicitly assumed, a small-velocity real spectral pair
produces ancient nonlinear orbits on the exact mass--momentum--energy
level and logarithmic orbital exit times. Finally, conserved internal
correlations and quantitative polarization give global energy-space
bounds beyond scalar critical mass for genuinely mixed states, and
force every unbounded critical-mass vector evolution to be exactly
scalar throughout its lifespan. Complete supporting proofs are included.
\end{abstract}

The traveling-wave nonlinear theorem retains one explicit real-sector
nondegeneracy assumption. The main amplification theorem does not
establish finite-time blowup, and the amplification exponent is not
claimed optimal. A primary-source comparison is given below.

\section*{The main connection}
The guiding question is how critical concentration geometry constrains
time evolution. For pure quartic dispersion, energy conservation alone
gives $A(t)\ge2|E|/d$. Our first result forces a substantially larger
kinetic scale, $c|E|d^{-5/3}$, somewhere along the trajectory. For data
$u_d(0)=(1+d)^{N/8}Q$, where $Q$ has mass $M_c$, this becomes
\[
 \sup_{0\le t<T_+}\frac{A(u_d(t))}{A(u_d(0))}\ge c d^{-2/3}.
\]
Thus for each prescribed factor $K$, sufficiently small nonzero $d$
forces a finite time at which the factor $K$ is exceeded. This result
concerns a time-dependent solution, without imposing its eventual form.

The exponent has a concrete source. Near-sharp interpolation controls
the part outside a constant complex phase by $O(\sqrt d)$. Local mass
conservation then bounds the speed of a concentration center by the
same order. Energy splitting leaves at most $O(d)$ mass in a remote
tail. Fourth-order interpolation charges a tail current at the rate
$L^{4/3}\mu$, producing $d^{2/3}d=d^{5/3}$. Annulus averaging is what
allows this argument to use arbitrary nonradial data.

The other two evolution results concern the same relationship from
different sides. Constrained phase adjustment determines how motion
changes a coherent wave; the spectral pair then yields genuine nonlinear
departure, with all conserved quantities fixed. In a unitary internal
system, polarization cannot be chosen anew at each time independently
of the conserved correlation operator. That conservation law turns a
static improvement of sharp interpolation into an all-time bound and
an exact scalar reduction at critical mass.

\paragraph{Organization.}
Section~\ref{sec:r10-foundations} fixes the conventions and compactness
input. Section~\ref{sec:r12-amplification} proves the nonradial kinetic
amplification result. The next sections treat constrained wave geometry,
nonlinear departure, and internal correlations. Appendices reproduce
the density geometry and spectral reduction used by these arguments.
The concrete strong $H^2$ flow proofs cover $N=1,2,3$; statements in
other dimensions or for general symbols retain their stated local-flow
assumptions.

\section{Analytic foundations and conventions}
\label{sec:r10-foundations}
The Fourier transform is unitary, $D=-i\nabla$, and
$P(u)=\operatorname{Im}\int\overline u\nabla u=\int\xi|\widehat u|^2$.
The real Hilbert pairing is the real part of the complex $L^2$ pairing.
Throughout the local-power arguments, $N\ge1$, $m>1$, $s=m/2$,
$r=m/N$, and $p=2+2r$. A principal symbol $a$ is real, even,
$m$-homogeneous, $C^{N+2}$ off zero, and satisfies
$c|\xi|^m\le a(\xi)\le C|\xi|^m$.
Set $A(u)=\int a|\widehat u|^2$, $M(u)=\|u\|_2^2$, and
$E(u)=A(u)/2-\|u\|_p^p/p$. The sharp critical mass $M_c$ is defined by
\begin{equation}\label{eq:r10-basic-GN}
 \|u\|_p^p\le\frac p2\left(\frac{M(u)}{M_c}\right)^r A(u).
\end{equation}
The traveling-wave convention $u(t,x)=e^{i\omega t}U(x-vt)$ gives
$[a(D)+\omega+iv\cdot\nabla]U-|U|^{2r}U=0$.

\begin{lemma}[Normalized optimizer compactness]
The inequality \eqref{eq:r10-basic-GN} is attained. If
$M(f_n)\to M_c$, $A(f_n)\to1$, and $\|f_n\|_p^p\to p/2$, then
some translated subsequence converges strongly in $H^s$ to an optimizer.
The normalized optimizer set is compact modulo translations and phases.
\end{lemma}
\begin{proof}
Sharp Sobolev interpolation gives a finite positive best constant, and
amplitude and spatial scaling normalize a maximizing sequence as stated.
Bounded energy-space sequences that vanish locally in mass converge to
zero in $L^p$: use a bounded-overlap cover by unit balls, local Sobolev
interpolation, and $2<p<2N/(N-m)$ when $N>m$ (all finite $p$ otherwise).
Thus translations give a nonzero weak limit $Q$ and almost-everywhere
convergence. Write $M_1=M(Q)>0$, $A_1=A(Q)>0$, $M_2=M_c-M_1$,
and $A_2=1-A_1$. Quadratic splitting and Brezis--Lieb give
\[
 M_c^r\le M_1^r A_1+M_2^r A_2,\qquad A_1+A_2=1.
\]
If both masses are positive, the right side is strictly below $M_c^r$.
Consequently $M_1=M_c$ and then $A_1=1$. Weak convergence plus
convergence of both quadratic norms gives strong $H^s$ convergence.
The same proof applies to a sequence of optimizers. Constant phase is
proved by the density-norm lemma below; this compactness argument itself
does not need a positivity-preserving kinetic operator.
\end{proof}

\begin{lemma}[Weak-profile regularity and Pohozaev identities]
Every weak $H^s$ local-power profile belongs to $L^\infty\cap H^{m+1}$
and satisfies
\begin{equation}\label{eq:r10-basic-pohozaev}
 v\cdot P=2mE,\qquad \omega M=rA+(2m-p)E.
\end{equation}
No spatial moment is required.
\end{lemma}
\begin{proof}
Choose $\Lambda$ so large that
$a(\xi)-v\cdot\xi+\Lambda\asymp1+|\xi|^m$.
The shifted resolvent maps $L^q$ to $H^{m,q}$, $1<q<\infty$,
by the Mihlin multiplier theorem. For $N>m$, the starting exponent
is $q_0=2N/(N-m)$. Splitting the nonlinearity into $|U|>1$ and
$|U|\le1$, resolvent Sobolev improvement iterates the reciprocal
exponent as $\beta_{j+1}=(1+2r)\beta_j-r$ with
$\beta_0=(1-r)/2$. The ideal iterates
$\beta_j=1/2-(r/2)(1+2r)^j$ cross the boundedness threshold after
finitely many steps. Inputs have exponent greater than one; intermediate
exponents may be weakened to avoid endpoints. The linear contribution
has the stronger improvement $\beta\mapsto\beta-r$ and is handled
separately. When $N=m$ start at any sufficiently large finite exponent;
when $N<m$ the initial space already embeds in $L^\infty$.
Now the nonlinearity is in $L^2$, giving $U\in H^m\subset H^1$.
The $C^1$ map $z\mapsto|z|^{2r}z$, with bounded derivative on the
range of $U$, sends $U$ into $H^1$; one more resolvent application
gives $H^{m+1}$.

Here are the cutoff details for the dilation calculation. Let
$f_R(x)=x\chi(x/R)$ with $\chi$ compactly supported and equal to
one near zero, and let $L_R=f_R\cdot\nabla+\operatorname{div}f_R/2$.
The Fourier kernel of $[a(D),L_R]$ is, up to the transform constant,
\[
 \frac i2(a(\xi)-a(\eta))(\xi+\eta)\cdot\widehat f_R(\xi-\eta).
\]
Splitting comparable and noncomparable frequencies gives
\[
 |a(\xi)-a(\eta)||\xi+\eta|
 \le C|\xi-\eta|(|\xi|+|\eta|)^m.
\]
After division by $\langle\xi\rangle^s\langle\eta\rangle^s$,
the convolution kernel is bounded by
$C|k|\langle k\rangle^s R^{N+1}|\widehat f(Rk)|$, with a
uniform $L^1$ norm for $R\ge1$. Hence the commutators are uniformly
bounded $H^s\to H^{-s}$. Schwartz approximation and the homogeneous
symbol identity give
$\langle U,[a(D),L_R]U\rangle\to mA(U)$.
The drift commutator tends to $-v\cdot P$; its order is one and
$s>1/2$. In the nonlinear pairing integration by parts and dominated
convergence give
\[
 \operatorname{Re}\int |U|^{p-2}\overline U L_RU
 \longrightarrow N(1/2-1/p)\|U\|_p^p=\frac m p\|U\|_p^p.
\]
Testing the profile with $L_RU$ therefore gives $mE=v\cdot P/2$.
Testing with $U$ gives $A+\omega M-v\cdot P-\|U\|_p^p=0$;
combining these two equations proves \eqref{eq:r10-basic-pohozaev}.
\end{proof}

The same regularity proof holds for an individually fixed nonnegative
order-$m$ addition $b$ with the stated Mihlin derivative bounds.
Constants need not be uniform in that addition or in $v$.
For nonhomogeneous minimizers, differentiation of their scalar
mass-preserving dilation energy is legitimate directly by dominated
convergence and minimality; differentiability of the dilation orbit
as an $H^s$-valued map is not being assumed.

\section{Forced kinetic amplification at nearly critical mass}
\label{sec:r12-amplification}

We return to the pure mass-critical equation
\[
 iu_t=\Delta^2u-|u|^{8/N}u,\qquad
 p=2+8/N,\qquad
 M=\|u\|_2^2,\quad A=\|\Delta u\|_2^2,\quad
 E=\frac A2-\frac{\|u\|_p^p}{p}.
\]
Let $M_c$ be the mass in the sharp interpolation inequality
\[
 \|f\|_p^p\le\frac p2\left(\frac{M(f)}{M_c}\right)^{4/N}A(f).
\]
Solutions are strong conservative $H^2$ solutions for which the integrated
localized mass and virial identities hold by approximation. No spatial
moment or radial symmetry is imposed.

\begin{theorem}[Quantitative kinetic amplification]
\label{thm:r12-amplification}
There are $d_0,C>0$, depending only on the scalar sharp interpolation
problem, with the following property. If a forward-global solution has
\[
 P=0,\qquad E=-e<0,\qquad
 d=\left(\frac M{M_c}\right)^{4/N}-1\in(0,d_0),\qquad
 B=\sup_{t\ge0}A(u(t))<\infty,
\]
then
\begin{equation}\label{eq:r12-amplification}
 e\le C d^{5/3}B.
\end{equation}
Consequently, in any maximal strong $H^2$ flow with the usual continuation
alternative,
\begin{equation}\label{eq:r12-maximal-amplification}
 \sup_{0\le t<T_+}A(u(t))\ge C^{-1}|E|d^{-5/3}.
\end{equation}
The latter statement allows finite-time blowup, global unbounded motion,
and bounded global motion.
\end{theorem}

The instantaneous sharp interpolation inequality gives only
$A(t)\ge2|E|/d$. Thus \eqref{eq:r12-maximal-amplification} forces a
larger kinetic scale, by a factor comparable to $d^{-2/3}$.
The result does not supply an upper bound on that scale and therefore
does not settle finite-time blowup or growup for all such initial data.

\subsection{A quantitative constant-phase estimate}
For a nonzero scalar complex field define
\[
 \mathfrak d(f)=1-
 \frac{\|f\|_p^p}{(p/2)(M(f)/M_c)^{4/N}A(f)}.
\]

\begin{lemma}[Distance to a real field]
\label{lem:r12-realphase}
There is $C<\infty$ such that every nonzero $f\in H^2(\mathbb R^N;\mathbb C)$
admits a constant phase $\theta$ for which
\begin{equation}\label{eq:r12-realphase}
 \frac{M(\operatorname{Im}(e^{-i\theta}f))}{M(f)}+
 \frac{A(\operatorname{Im}(e^{-i\theta}f))}{A(f)}
 \le C\mathfrak d(f).
\end{equation}
No uniqueness of real optimizers or nondegeneracy of the real linearized
operator is assumed.
\end{lemma}
\begin{proof}
Amplitude and critical spatial scaling reduce the proof to $M(f)=M_c$
and $A(f)=1$. If the result failed, since the left side is at most two,
there would be a sequence with $\mathfrak d(f_n)\to0$ for which the ratio
of the right defect to the best left side tends to zero. Sharp-optimizer
compactness supplies, after translations and phases,
$f_n\to Q$ strongly in $H^2$, for a real normalized optimizer $Q$.
Choose the remaining phase so that $f_n=X_n+iY_n$ and
$\langle Q,Y_n\rangle=0$.
Put $r=4/N$, $\kappa=p/2=1+r$, $a_n=A(Y_n)$ and $y_n=M(Y_n)$.
The transverse Taylor expansion, with the Sobolev remainder estimate,
gives
\[
 \|X_n+iY_n\|_p^p
 =\|X_n\|_p^p+\kappa\int |X_n|^{p-2}Y_n^2
       +o(\|Y_n\|_{H^2}^2).
\]
For $2<p\le4$ the remainder is bounded by $C\|Y_n\|_p^p$;
for $p>4$ use the sum of the fourth-order mixed term and
$C\|Y_n\|_p^p$. Sharp interpolation for $X_n$ yields
\[
 \|X_n\|_p^p\le\kappa(1-y_n/M_c)^r(1-a_n).
\]
Potential form convergence therefore implies
\[
 \kappa-\|f_n\|_p^p
 \ge\kappa\langle Y_n,L_-Y_n\rangle-o(\|Y_n\|_{H^2}^2),
 \qquad L_-=\Delta^2+r/M_c-|Q|^{p-2}.
\]
The sharp phase form has kernel exactly $\mathbb RQ$ and is coercive
on $Q^\perp$, as follows from density strict convexity and the sharp
interpolation inequality. Hence the defect controls
$\|Y_n\|_{H^2}^2$, a contradiction. Compactness of normalized optimizers
makes the coercivity constant uniform. This proves the lemma.
\end{proof}

This is a real two-component polarization statement. Treating $f$ as a
single complex internal component would give a vacuous rank-one assertion
and would not prove \eqref{eq:r12-realphase}.

\subsection{A slowly moving concentration center}
Critical scaling allows us to assume $B=1$. Throughout the proof $M$ then
lies in a fixed compact interval about $M_c$, and
\begin{equation}\label{eq:r12-Arange}
 0<2e/d\le A(t)\le1.
\end{equation}
Since
\[
 \mathfrak d(u(t))=\frac{d-2e/A(t)}{1+d}\le d,
\]
Lemma~\ref{lem:r12-realphase} gives, separately at each time, a phase
such that $u=e^{i\theta}(X+iY)$ with
\begin{equation}\label{eq:r12-phase-small}
 \|Y\|_{H^2}\le C\sqrt d,\qquad \|X\|_{H^2}\le C.
\end{equation}
No time derivative of the phase is used.

For a real smooth bounded test function $\phi$, integration by parts in
the mass identity gives
\begin{equation}\label{eq:r12-density-id}
 \frac d{dt}\int\phi|u|^2
 =2\operatorname{Im}\int
 \bigl(2\nabla\phi\cdot\nabla\overline u+(\Delta\phi)\overline u\bigr)
             \Delta u.
\end{equation}
The local nonlinearity cancels. Substituting
\eqref{eq:r12-phase-small} proves
\[
 \left|\frac d{dt}\int\phi|u|^2\right|
 \le C\sqrt d\bigl(\|\nabla\phi\|_\infty+\|\Delta\phi\|_\infty\bigr).
\]
In particular, for $R\ge1$ and fixed $b$,
\begin{equation}\label{eq:r12-density-speed}
 \left|\frac d{dt}\int\chi((x-b)/R)|u|^2\right|
 \le C\sqrt d/R.
\end{equation}
These statements hold in integrated form for the stated solution class.

Choose $\varepsilon_0=M_c/32$. For sufficiently small $d_0$, compactness
of sharp near-optimizers and \eqref{eq:r12-Arange} give a radius
$R_0\ge1$ and a choice $c(t)$ satisfying
\begin{equation}\label{eq:r12-core-mass}
 \int_{|x-c(t)|\le R_0}|u(t,x)|^2\,dx\ge M_c-\varepsilon_0.
\end{equation}
The radius may depend on $e/d$; no uniform lower kinetic scale is needed.
Take $\chi=1$ on the unit ball and $\chi=0$ outside the radius-two ball.
If $|c(t)-c(s)|>4R_0$, the integral in
\eqref{eq:r12-density-speed}, centered at $c(s)$, changes by at least
$M_c-3\varepsilon_0$, provided $M-M_c\le\varepsilon_0$.
Thus, for a fixed $\tau>0$,
\[
 |t-s|\le\tau R_0/\sqrt d\quad\Longrightarrow\quad
 |c(t)-c(s)|\le4R_0.
\]
Interpolate $c$ linearly at time spacing $\tau R_0/\sqrt d$. We obtain
one Lipschitz center $a$ satisfying
\begin{equation}\label{eq:r12-center-speed}
 \|a'\|_\infty\le L=C\sqrt d,\qquad
 \sup_t|a(t)-c(t)|\le8R_0.
\end{equation}
Crucially, its speed bound is independent of $R_0$ and $e/d$.
A fixed ball about $a(t)$ retains the core mass in
\eqref{eq:r12-core-mass}.

\begin{lemma}[Uniform mass outside the core]
\label{lem:r12-tailmass}
For every $\eta>0$ there is $R_\eta<\infty$ such that
\begin{equation}\label{eq:r12-tailmass}
 \sup_{t\ge0}\int_{|x-a(t)|>R_\eta}|u(t,x)|^2\,dx
 \le M-M_c+\eta.
\end{equation}
\end{lemma}
\begin{proof}
Otherwise choose $t_k$ and $R_k\to\infty$ violating the inequality.
After translation by $a(t_k)$, a subsequence converges weakly in $H^2$
and strongly locally in $L^2$ to $v$. The fixed core gives
$M(v)\ge M_c-\varepsilon_0$, while the assumed exterior mass implies
$M(v)\le M_c-\eta$. The remainder $w_k$ therefore has limiting mass
at most $M-M_c+\varepsilon_0<M_c$. Quadratic splitting and the
Brezis--Lieb identity give
\[
 -e=E(v)+E(w_k)+o(1).
\]
Both terms on the right are nonnegative by sharp interpolation, since
both masses are strictly subcritical. This is a contradiction.
\end{proof}

The lemma permits a noncompact mass tail of size $M-M_c$. It establishes
neither mass compactness nor kinetic compactness of the full trajectory.

\subsection{Negative exterior kinetic energy in the virial derivative}
Choose a real radial $\phi_R$ equal to $|x|^2/2$ on $|x|\le R$ and
constant on $|x|\ge2R$, with
\[
 D^2\phi_R\le I,\qquad
 \|D^j\phi_R\|_\infty\le C_jR^{2-j}.
\]
One construction prescribes $\phi_R'(r)=rh(r/R)$ with smooth
nonincreasing $h$, equal to one on $[0,1]$ and zero on $[2,\infty)$.
Set
\begin{align*}
 J_R(t)&=2\operatorname{Im}\int
 \nabla\phi_R(x-a(t))\cdot\overline u\nabla u,\\
 K_R(t)&=\int_{|x-a(t)|>2R}|D^2u|^2,\\
 H_R(t)&=\int_{R/2<|x-a(t)|<2R}|D^2u|^2.
\end{align*}
We use the full Hessian in these tails. Its total squared norm equals
$\|\Delta u\|_2^2$, but the two densities are not pointwise identical.

The exact localized quartic identity is
\begin{align*}
 J_R'={}&8\operatorname{Re}\int\partial_{jk}\phi_R\,
           \partial_{ij}\overline u\,\partial_{ik}u
 -4\operatorname{Re}\int\partial_{jk}\Delta\phi_R\,
           \partial_j\overline u\,\partial_ku\\
 &-2\int\Delta^2\phi_R|\nabla u|^2
   +\int\Delta^3\phi_R|u|^2
   -\frac{16}{Np}\int\Delta\phi_R|u|^p
   -2a'\cdot\int D^2\phi_R\operatorname{Im}(\overline u\nabla u).
\end{align*}
The cutoff is centered at $a(t)$. Relative to the global kinetic term
$8A$, the first term has a nonpositive defect everywhere, and its defect
is exactly $-8|D^2u|^2$ outside $2R$. Higher cutoff derivatives contribute
$O(R^{-2})$. Subtracting the total momentum from the last term and using
$P=0$ therefore gives
\begin{equation}\label{eq:r12-retained-virial}
 J_R'\le-16e-8K_R+C\int_{|x-a|>R}|u|^p
          +CR^{-2}+CL\int_{|x-a|>R}|u||\nabla u|.
\end{equation}
Also $|J_R|\le CR$.

Fix small $\eta>0$, put $\mu=M-M_c+\eta$, and take $R\ge2R_\eta$.
For a cutoff $\zeta_R$ vanishing inside $R/2$ and equal to one outside
$R$, Lemma~\ref{lem:r12-tailmass} and the product rule give
\[
 M(\zeta_Ru)\le\mu,\qquad
 A(\zeta_Ru)\le C(K_R+H_R+R^{-2}).
\]
Sharp interpolation and
$\|\nabla f\|_2^2\le\|f\|_2\|\Delta f\|_2$ imply
\begin{align*}
 \int_{|x-a|>R}|u|^p
   &\le C\mu^{4/N}(K_R+H_R+R^{-2}),\\
 \int_{|x-a|>R}|u||\nabla u|
   &\le C\mu^{3/4}\bigl((K_R+H_R)^{1/4}+R^{-1/2}\bigr).
\end{align*}
Choose $d_0$ and $\eta$ so small that $C\mu^{4/N}\le1$.
Absorb the terms containing $K_R$ into its negative coefficient in
\eqref{eq:r12-retained-virial}; for the current term use Young's inequality.
We obtain
\begin{equation}\label{eq:r12-annular-virial}
 J_R'\le-16e+C\mu^{4/N}H_R+CL\mu^{3/4}H_R^{1/4}
             +CL^{4/3}\mu+o_R(1),
\end{equation}
where $o_R(1)\to0$ uniformly in time. No favorable lower bound on the
kinetic defect in the transition annulus was used.

\subsection{Averaging disjoint annuli}
Fix a positive integer $K$ and $R_*\ge2R_\eta$. Let
$R_j=8^jR_*$ for $0\le j<K$. Their transition annuli are disjoint about
the same center, so
\[
 \sum_{j=0}^{K-1}H_{R_j}(t)\le A(t)\le1.
\]
For each $T$ choose an index whose time-averaged $H_{R_j}$ is at most
$K^{-1}$. Integrate \eqref{eq:r12-annular-virial}, use Jensen's inequality
for the fourth root, and bound both endpoint values of $J_{R_j}$:
\[
 16e\le\frac{2C8^{K-1}R_*}{T}
       +\frac{C\mu^{4/N}}K+CL\mu^{3/4}K^{-1/4}
       +CL^{4/3}\mu+\sup_{R\ge R_*}|o_R(1)|.
\]
First let $T\to\infty$ with $K,R_*$ fixed, then $K\to\infty$,
then $R_*\to\infty$, and finally $\eta\downarrow0$. This proves
\[
 e\le CL^{4/3}(M-M_c).
\]
The radius is fixed during each time integration. By
\eqref{eq:r12-center-speed}, $L\le C\sqrt d$, and $M-M_c\le Cd$.
Thus $e\le Cd^{5/3}$ when $B=1$. Undoing critical scaling proves
Theorem~\ref{thm:r12-amplification}.

\subsection{Finite-time excursions from every sharp standing wave}
\begin{corollary}[Standing-wave instability without real-sector nondegeneracy]
\label{cor:r12-standing-instability}
Let $Q$ be any real sharp optimizer normalized by $M(Q)=M_c$,
with its associated standing-wave frequency. Assume a strong $H^2$ local flow and the continuation alternative.
For
\[
 u_d(0)=(1+d)^{N/8}Q,\qquad d>0,
\]
and every prescribed $K>1$, all sufficiently small $d$ admit a finite
forward time $t_d<T_+(d)$ for which
\[
 A(u_d(t_d))>K A(u_d(0)).
\]
In particular the standing wave is orbitally unstable in $H^2$ modulo
constant phase and translations. No radial or $L_+$ nondegeneracy
hypothesis is needed.
\end{corollary}
\begin{proof}
These initial data have $P=0$, $M=(1+d)^{N/4}M_c$ and
$E=-dA(u_d(0))/2$. Theorem~\ref{thm:r12-amplification} gives
\[
 \sup_{0\le t<T_+}\frac{A(u_d(t))}{A(u_d(0))}
 \ge(2C)^{-1}d^{-2/3}.
\]
For small $d$ this is strictly greater than $K$, so a finite time has
the claimed property. Since $u_d(0)\to Q$ in $H^2$, take, for example,
$K=4$ and use
\[
 \inf_{\theta,y}\|u_d(t_d)-e^{i\theta}Q(\cdot-y)\|_{H^2}
 \ge\|\Delta u_d(t_d)\|_2-\|\Delta Q\|_2
 \ge\|\Delta Q\|_2.
\]
This proves orbital instability.
\end{proof}

For $N=1,2,3$ the required strong local flow follows from the $H^2$
algebra and composition estimates, with conservation and the integrated
identities obtained by smooth approximation. In larger dimensions the statement
uses the appropriate strong energy-space local theory. The conclusion
is a finite-time excursion for the perturbed data; its proof does not
identify which of finite blowup, global unbounded evolution, or bounded
large motion occurs.

\section{The energy geometry at fixed mass and momentum}
\label{sec:r12-constraint}

The spectral reduction for a specified smooth branch requires control of
the real linearized kernel. The variational assertion below does not:
the imaginary response alone forces a negative constrained energy
curvature when the homogeneous order exceeds two. Every component of
momentum is preserved, and no differentiability of the dilation group on
the individual profile is assumed.

Let $s=m/2$, $\alpha=2m/N$, and let $a$ be a positive, real-even,
elliptic homogeneous symbol of degree $m>1$. Write
\[
 E(u)=\frac12 A(u)-\frac1{\alpha+2}\|u\|_{\alpha+2}^{\alpha+2},
 \quad M(u)=\|u\|_2^2,
 \quad P(u)=\operatorname{Im}\int\overline u\nabla u.
\]
Let $Q$ be a real sharp ground profile at frequency one. The scalar
phase-coercivity theorem gives
\[
 L_-=a(D)+1-|Q|^\alpha\geq0,\qquad \ker L_-=\mathbb RQ,
\]
with coercivity on $H^s\cap Q^\perp$. Define
\begin{equation}\label{eq:r12-phase-T}
 Z_j=L_-^{-1}\partial_jQ,\qquad
 T_{ij}=\langle\partial_iQ,Z_j\rangle.
\end{equation}
The inverse and pairing can be understood in energy-space form sense.
In particular $T$ is positive definite.

\begin{theorem}[Constrained curvature without a profile branch]
\label{thm:r12-constrained-curvature}
Suppose $v_n\to0$, $v_n\ne0$, and exact profiles satisfy
\[
 (a(D)+1+iv_n\cdot\nabla)U_n-|U_n|^\alpha U_n=0,
 \qquad U_n\longrightarrow Q\quad\hbox{strongly in }H^s,
\]
after phase and translation. For every sufficiently large $n$ there is
a continuous $H^s$ curve $V_n(\lambda)$, defined near $\lambda=1$, with
\[
 V_n(1)=U_n,\qquad
 M(V_n(\lambda))=M(U_n),\qquad
 P(V_n(\lambda))=P(U_n).
\]
Its scalar energy is $C^2$ and obeys
\begin{equation}\label{eq:r12-constrained-curvature}
 \left.\frac{d}{d\lambda}E(V_n(\lambda))\right|_1=0,
 \qquad
 \left.\frac{d^2}{d\lambda^2}E(V_n(\lambda))\right|_1
 =(2-m)v_n^TTv_n+o(|v_n|^2).
\end{equation}
For $m>2$ the profile is therefore not a local energy minimum at its
fixed mass and momentum: both sides of this curve have strictly lower
energy. The statement assumes neither a differentiable family of
profiles nor nondegeneracy of the real block $L_+$.
\end{theorem}

\begin{proof}
We suppress the sequence index and keep errors relative to its limit.
Choose the phase so that $U=X+iY$ and $\langle Y,Q\rangle=0$.
The imaginary profile equation is
\[
 L_-Y=-v\cdot\nabla X+(|U|^\alpha-|Q|^\alpha)Y.
\]
Put $p=\alpha+2$. Since $H^s\hookrightarrow L^p$, the potential
$|U|^\alpha-|Q|^\alpha$ tends to zero as a quadratic form on $H^s$:
this follows from H\"older and its convergence in $L^{p/\alpha}$.
The derivative maps $H^s$ continuously to $H^{-s}$ because $2s=m>1$.
Project onto $Q^\perp$ and apply phase coercivity. First
$\|Y\|_{H^s}=O(|v|)$; substituting this bound back into the equation gives
\begin{equation}\label{eq:r12-imaginary-response}
 Y=-\sum_jv_jZ_j+o(|v|)\quad\hbox{in }H^s.
\end{equation}
No estimate for a derivative of the real profile is needed.
The bilinear momentum form is continuous on $H^s$, and
$P_j(U)=-2\langle Y,\partial_jX\rangle$. Hence
\begin{equation}\label{eq:r12-P-response}
 P(U)=2Tv+o(|v|).
\end{equation}
The positivity of $T$ follows from phase coercivity and the fact that a
nonzero $L^2$ function cannot have a vanishing nonzero directional
derivative.

For $c=(c_0,c_1,\ldots,c_N)\in\mathbb R^{N+1}$ set
\[
 W(c)=(1+c_0)U+i\sum_{j=1}^N c_jZ_j.
\]
The Jacobian at $c=0$ of $c\mapsto(M(W(c)),P(W(c)))$ is
\[
 \begin{pmatrix}
  2M(U)&O(|v|)\\
  2P(U)&-2T+o(1)
 \end{pmatrix}.
\]
It is invertible for sufficiently small velocity. The finite-dimensional
implicit-function theorem supplies a smooth $c(\lambda)$ with $c(1)=0$
and
\begin{equation}\label{eq:r12-W-charges}
 M(W(\lambda))=M(U),\qquad
 P(W(\lambda))=\lambda^{-1}P(U).
\end{equation}
Here and below $W(\lambda)$ denotes $W(c(\lambda))$.
Differentiating these equations and using
\eqref{eq:r12-P-response} gives
\[
 c_j'(1)=v_j+o(|v|),\qquad c_0'(1)=O(|v|^2),
 \qquad h:=W'(1)=i\sum_jv_jZ_j+o(|v|)\quad\hbox{in }H^s.
\]
The amplitude estimate uses
$\langle U,iZ_j\rangle=\langle Y,Z_j\rangle=O(|v|)$.

Define
\[
 V(\lambda,x)=\lambda^{N/2}W(\lambda,\lambda x).
\]
Critical homogeneity gives the exact identities
\begin{equation}\label{eq:r12-scaled-charges}
 M(V(\lambda))=M(U),\quad P(V(\lambda))=P(U),\quad
 E(V(\lambda))=\lambda^mE(W(\lambda)).
\end{equation}
Strong continuity of dilation gives a continuous $H^s$ curve. Although
that curve need not be differentiable in $H^s$, its scalar energy is
$C^2$ by the last identity. This is the reason that no weighted moment
or dilation-domain assumption enters the theorem.

Let $\mathcal L_U$ be the real Hessian of
$S=E+M/2-v\cdot P/2$, and put
$q_U(h)=\langle h,\mathcal L_Uh\rangle$ in form notation.
Because $S'(U)=0$, differentiating \eqref{eq:r12-W-charges} yields
\[
 \left.\frac{d}{d\lambda}E(W(\lambda))\right|_1=-\frac12v\cdot P(U),
 \qquad
 \left.\frac{d^2}{d\lambda^2}E(W(\lambda))\right|_1
 =q_U(h)+v\cdot P(U).
\]
Using the exact profile identity $v\cdot P=2mE$ in
\eqref{eq:r12-scaled-charges}, we conclude that
\begin{equation}\label{eq:r12-exact-curvature}
 E(V)'(1)=0,\qquad
 E(V)''(1)=q_U(h)-\frac{m-1}{2}v\cdot P(U).
\end{equation}
The energy is a $C^2$ functional on $H^s$, since $p>2$ and
$H^s\hookrightarrow L^p$. Its Hessian forms therefore converge on bounded
sets as $U\to Q$, $v\to0$. The imaginary block at $Q$ is $L_-$, so
\[
 q_U(h)=q_Q\Big(i\sum_jv_jZ_j\Big)+o(|v|^2)
 =v^TTv+o(|v|^2).
\]
Together with \eqref{eq:r12-P-response} and
\eqref{eq:r12-exact-curvature}, this proves
\eqref{eq:r12-constrained-curvature}. The sign for $m>2$ follows from
positive definiteness of $T$.
\end{proof}

\begin{corollary}[All sufficiently slow action ground waves]
\label{cor:r12-all-ground-waves}
At frequency one, define
\[
 K_v(u)=A(u)+M(u)-v\cdot P(u),\qquad
 \kappa_v=\inf_{u\ne0}\frac{K_v(u)}{\|u\|_p^2},
 \qquad p=2+\frac{2m}{N}.
\]
For sufficiently small $|v|$, this infimum is attained. Normalize each
minimizer $U$ by $K_v(U)=\|U\|_p^p$; these are precisely the action
minimizers on the nonzero Nehari manifold. If $m>2$, then every such
ground wave at every sufficiently small nonzero velocity is not a local
energy minimum at fixed mass and momentum. This conclusion requires no
uniqueness, differentiable branch, or real-profile nondegeneracy.
\end{corollary}

\begin{proof}
Ellipticity and H\"older in Fourier variables give
\[
 |P(u)|\leq C A(u)^{1/m}M(u)^{1-1/m}
          \leq C'(A(u)+M(u)).
\]
Consequently $K_v$ is uniformly equivalent to the $H^s$ norm squared
for sufficiently small $|v|$. The exponent $p>2$ is strictly below the
energy-space Sobolev endpoint, including when $s\geq N/2$.
For completeness, normalize a minimizing sequence by $\|u_n\|_p=1$.
It is bounded in $H^s$ and cannot vanish locally, by the standard
localized Sobolev interpolation argument. After a translation it has a
nonzero weak limit $u$. Local compact embedding and the
Brezis--Lieb identity split the $p$th power norm into $c$ and $1-c$,
where $c=\|u\|_p^p>0$. Weak Hilbert-space orthogonality splits $K_v$.
Thus
\[
 \kappa_v\geq\kappa_v\big(c^{2/p}+(1-c)^{2/p}\big).
\]
Since $2/p<1$, this forces $c=1$. Norm convergence and the coercivity of
$K_v$ give strong $H^s$ convergence and attainment. Rescaling a minimizer
to the Nehari manifold gives the profile equation and
\begin{equation}\label{eq:r12-ground-norm}
 K_v(U)=\|U\|_p^p=\kappa_v^{p/(p-2)}.
\end{equation}
On that manifold $S_v(U)=(1/2-1/p)\|U\|_p^p$, proving the asserted
minimizing equivalence.

The Fourier bound also gives $\kappa_v\to\kappa_0$ as $v\to0$.
Let $v_n\to0$ and let $U_n$ be any Nehari-normalized minimizers.
They are bounded in $H^s$ by \eqref{eq:r12-ground-norm}, and
$K_0(U_n)-K_{v_n}(U_n)=v_n\cdot P(U_n)\to0$. Therefore their normalized
versions form a minimizing sequence for $\kappa_0$.
The preceding compactness proof gives, after translations and passing to
a subsequence, $U_n\to Q$ strongly in $H^s$, where $Q$ is a
frequency-one ground profile.

We verify that this limiting profile is one of the sharp critical
optimizers used in Theorem~\ref{thm:r12-constrained-curvature}.
Write $r=m/N$ and let the sharp critical inequality be
\[
 \|u\|_p^p\leq C_* M(u)^r A(u).
\]
For fixed $A+M$, the product $M^rA$ is maximal at $M=rA$.
A critical dilation of a sharp optimizer attains that ratio, so it also
attains $\kappa_0$. Conversely, attainment of $\kappa_0$ forces equality
in both inequalities, hence sharp critical optimization and $M=rA$.
The density phase-rigidity theorem makes this optimizer real after a
constant phase. Nehari normalization then gives
$\|Q\|_p^p=(1+r)A(Q)$, and equality in
$\|Q\|_p^p\leq(1+r)(M(Q)/M_c)^r A(Q)$ gives $M(Q)=M_c$.

If the asserted uniform small-velocity conclusion failed, there would
be a sequence of nonzero $v_n\to0$ and ground waves lacking a negative
fixed-charge curve. The compactness just proved produces a real limiting
$Q$, to which Theorem~\ref{thm:r12-constrained-curvature} applies. Its
strictly negative curvature contradicts that failure. No uniform
choice of a differentiable family was used.
\end{proof}

\paragraph{Why the phase correction matters.}
The linear boost
\[
 V_b(\lambda,x)=e^{i(1-\lambda)k\cdot x}
                  \lambda^{N/2}U(\lambda x),\qquad
 k=\frac{P(U)}{M(U)},
\]
also fixes both charges. However its curvature, when $a$ has the
homogeneous $C^2$ derivative bounds and $m>2$, is
\begin{equation}\label{eq:r12-linear-boost}
 E(V_b)''(1)=\frac12\left(k^TH_Uk-(m-1)v\cdot P(U)\right),
 \qquad H_U=\int D^2a(\xi)|\widehat U(\xi)|^2\,d\xi.
\end{equation}
Indeed the kinetic term is
$\int a(k+\lambda(\xi-k))|\widehat U(\xi)|^2\,d\xi$.
Euler's identities and the profile boost identity
$\int\nabla a|\widehat U|^2=vM$ give
\eqref{eq:r12-linear-boost}. Near $Q$ its leading term is
\[
 \frac{2}{M(Q)^2}v^TTH_QTv-(m-1)v^TTv.
\]
Phase positivity alone gives no sign for this expression.
By contrast the optimal imaginary correction minimizes
$\langle\psi,L_-\psi\rangle$ under
$-2\langle\partial_jQ,\psi\rangle=-P_j$. Its minimum is
$P^TT^{-1}P/4$, attained at
$\psi=\frac12\sum_j(T^{-1}P)_jZ_j$.
This is the correction used to leading order in the theorem.

\begin{remark}[Energetic and dynamical conclusions]
The theorem gives energy-decreasing perturbations at precisely the same
conserved charges as the wave. It does not alone prove that their
solutions escape a neighborhood of the symmetry orbit. In particular,
Morse index one does not remove an uncontrolled extra kernel: a negative
direction can belong to a zero Jordan chain. Nonlinear or spectral
instability requires the additional dynamical argument stated in its
own theorem. For $1<m<2$ the displayed curve instead has positive
leading curvature, which also does not establish stability.
\end{remark}

\section{Nonlinear departure on the exact invariant level}
\label{sec:r12-nonlinear}

We now turn the real spectral pair into an actual nonlinear trajectory.
The equation throughout this section is the pure mass-critical quartic
model
\begin{equation}\label{eq:r12-pure-quartic}
 iu_t=\Delta^2u-|u|^{8/N}u,\qquad N\in\{1,2,3\},
 \qquad X=H^2(\mathbb R^N;\mathbb C).
\end{equation}
The space $X$ is regarded as real. In this concrete phase space, the
branch, weighted regularity, and complementary gap used in the spectral theorem in Appendix~\ref{sec:r11-spectral} follow from real-profile nondegeneracy. That remaining
assumption is stated explicitly below. The form decomposition is the
one-negative-direction instance of the Hamiltonian trichotomy developed
by Lin and Zeng~\cite{LZ}; the estimates needed here are proved directly.

\begin{theorem}[Ancient departure and nonlinear orbital instability]
\label{thm:r12-nonlinear-instability}
Let $Q$ be a real sharp ground profile at frequency one. Assume
\begin{equation}\label{eq:r12-real-nondegeneracy}
 \ker\left(\Delta^2+1-\left(1+\frac8N\right)|Q|^{8/N}\right)
   =\operatorname{span}_{\mathbb R}
        \{\partial_1Q,\ldots,\partial_NQ\}.
\end{equation}
Then there is a local $C^2$ profile family $U_{\omega,v}$ through $Q$,
unique up to phase and translation near $Q$. Put $U_v=U_{1,v}$.
For every sufficiently small fixed $v\ne0$, put
\[
 \mathcal O_v=\{e^{i\theta}U_v(\cdot-y):\theta\in\mathbb R,
                                     \ y\in\mathbb R^N\}.
\]
There is a nonstationary solution in the co-moving, co-rotating frame,
defined for $t\le0$, of the form
\begin{equation}\label{eq:r12-ancient-asymptotic}
 w_v(t)=U_v+\varepsilon_v e^{\lambda_vt}Y_{+,v}
              +O_v(e^{(1+\beta)\lambda_vt}),\qquad t\to-\infty,
\end{equation}
where $\|Y_{+,v}\|_{H^2}=1$, $J\mathcal L_vY_{+,v}
=\lambda_vY_{+,v}$, $\varepsilon_v>0$, and
\[
 \beta=\begin{cases}1,&N=1,2,\\2/3,&N=3.\end{cases}
\]
At every time this solution has exactly the same mass, momentum, and
energy as $U_v$.

Consequently, there are $\rho_v>0$ and data $u_{\delta,v}$ satisfying
\begin{gather*}
 M(u_{\delta,v})=M(U_v),\quad P(u_{\delta,v})=P(U_v),\quad
 E(u_{\delta,v})=E(U_v),\\
 u_{\delta,v}=U_v+\delta Y_{+,v}+O_v(\delta^{1+\beta})
                         \quad\hbox{in }H^2,
\end{gather*}
whose first exit time from the orbital distance $\rho_v$ obeys
\begin{equation}\label{eq:r12-exit-time}
 \tau_{\delta,v}=\lambda_v^{-1}\log(1/\delta)+O_v(1)
 \qquad(\delta\downarrow0).
\end{equation}
The spectral calculation gives
\begin{equation}\label{eq:r12-rate}
 \lambda_v=\sqrt{\frac{2v^TTv}{\mathcal D_Q}}+o(|v|),
\end{equation}
where
\begin{gather*}
 T_{ij}=\langle\partial_iQ,L_-^{-1}\partial_jQ\rangle,
 \qquad c_j=\langle\Lambda Q,L_-^{-1}\partial_jQ\rangle,\\
 \mathcal D_Q=\langle\Lambda Q,L_-^{-1}\Lambda Q\rangle-c^TT^{-1}c.
\end{gather*}
In particular, $e^{it}U_v(x-vt)$ is orbitally unstable in $H^2$ even
on its exact joint mass--momentum--energy level.
\end{theorem}

\subsection{Local flow and the nonlinear remainder}

Set $F(z)=|z|^{8/N}z$, with $\mathbb C$ identified with $\mathbb R^2$.
Since $H^2\hookrightarrow L^\infty\cap W^{1,4}$ for $N\le3$ and
$H^2$ is an algebra, $F:X\to X$ is locally $C^{1,\beta}$ with the
exponent in the theorem. Here is a direct check of the regularity that
is needed. Twice differentiating $DF(w)h$ produces sums of
\begin{gather*}
 DF(w)\nabla^2h,\qquad D^2F(w)[\nabla w,\nabla h],\\
 D^2F(w)[\nabla^2w,h],\qquad D^3F(w)[\nabla w,\nabla w,h].
\end{gather*}
The products lie in $L^2$ by the stated embeddings. On bounded subsets
of $\mathbb R^2$, $D^3F$ is Lipschitz in dimensions one and two, and
$2/3$-H\"older in dimension three. Subtracting these expressions at
$w_1,w_2$ therefore gives
\[
 \|DF(w_1)-DF(w_2)\|_{\mathcal B(X)}
       \le C_R\|w_1-w_2\|_X^\beta
\]
for $\|w_i\|_X\le R$ and $\|w_1-w_2\|_X\le1$.
The lower derivative terms obey the same bound. Smooth approximation
justifies the weak chain rules. It follows that, near a fixed $U$,
\begin{align}
 \|F(U+h)-F(U)-DF(U)h\|_X&\le C\|h\|_X^{1+\beta},
 \label{eq:r12-remainder}\\
 \|R(h_1)-R(h_2)\|_X
 &\le C(\|h_1\|_X^\beta+\|h_2\|_X^\beta)\|h_1-h_2\|_X,
 \label{eq:r12-remainder-lipschitz}
\end{align}
where $R(h)=i[F(U+h)-F(U)-DF(U)h]$.

The free group $e^{-it\Delta^2}$ is unitary on $X$. A contraction
argument gives local well-posedness and continuation while the $H^2$
norm is bounded. Mass, momentum, and energy are conserved. For a direct
justification, insert a selfadjoint Fourier cutoff $P_n$ into the
Hamiltonian nonlinearity $P_nF(P_nu)$, differentiate the conserved
quantities for this regularized flow, and pass to the $H^2$ solution
uniformly on compact existence intervals. The cutoff commutes with
phase and translation actions. Thus no Strichartz or spectral-mapping
hypothesis is required for the local nonlinear flow.

For
\[
 w(t,x)=e^{-it}u(t,x+vt)=U_v+h(t,x)
\]
the perturbation equation is
\begin{equation}\label{eq:r12-perturbation-eq}
 h_t=\mathcal A_vh+R_v(h),\qquad
 \mathcal A_v=J\mathcal L_v,\quad J=-i.
\end{equation}
The phase and translation transformations are $H^2$ isometries and
preserve all three conserved quantities.

\subsection{Why the other spectral assumptions follow here}

\begin{lemma}[Local branch and complementary gap]
\label{lem:r12-branch}
For a real sharp quartic ground profile satisfying
\eqref{eq:r12-real-nondegeneracy}, all the other assumptions of the
small-velocity theorem in Appendix~\ref{sec:r11-spectral} hold. In particular,
$\Lambda Q\in H^4$, the local profile branch is $C^2$ into $H^4$,
and its only Hessian zero modes are phase and translations.
\end{lemma}
\begin{proof}
Sharp phase coercivity gives $L_-\ge0$, $\ker L_-=\mathbb RQ$,
and an inverse on $Q^\perp$. The nonlinearity calculation above and
$(\Delta^2+1)Q=F(Q)$ give $Q\in H^6$. Therefore
$V=|Q|^{8/N}\in W^{2,\infty}$ and $V(x)\to0$ as $|x|\to\infty$.
We first justify the scaling derivative without assuming a moment.

Choose smooth bounded truncations $\phi_R$ of
$a\langle x\rangle$, with $\phi_R\to a\langle x\rangle$ pointwise,
$\phi_R\le Ca\langle x\rangle$, and
$|\nabla\phi_R|+|\Delta\phi_R|\le Ca$ uniformly in $R$.
Test $(\Delta^2+1-V)Q=0$ with $e^{2\phi_R}Q$ and put
$w_R=e^{\phi_R}Q$. The exact real identity
\[
 \operatorname{Re}\langle\Delta^2Q,e^{2\phi_R}Q\rangle
 =\|\Delta w_R+|\nabla\phi_R|^2w_R\|_2^2
  -\|2\nabla\phi_R\cdot\nabla w_R+(\Delta\phi_R)w_R\|_2^2
\]
shows that the difference from $\|\Delta w_R\|_2^2$ is bounded by
$Ca^2\|w_R\|_{H^2}^2$ for $a\le1$.
Choose a fixed ball outside which $V\le1/4$. Inside the ball,
$\int V|w_R|^2$ is uniformly bounded because the weights there are
bounded independently of $R$. Taking $a>0$ small absorbs both the
exterior potential and the $Ca^2$ error. Thus $w_R$ is uniformly
bounded in $H^2$, and passage to the limit gives
$w=e^{a\langle x\rangle}Q\in H^2$.

The conjugated equation
\[
 e^{a\langle x\rangle}(\Delta^2+1)e^{-a\langle x\rangle}w=Vw
\]
has smooth bounded coefficients in its lower-order terms, of order at
most three. Since $V\in W^{2,\infty}$, elementary elliptic bootstrapping
improves $w$ successively from $H^2$ to $H^3,H^4,H^5,H^6$.
At the first step, the third-order terms belong to $H^{-1}$; at the
subsequent steps they belong to $L^2,H^1,H^2$, respectively.
Consequently $\Lambda Q=NQ/2+x\cdot\nabla Q\in H^4$, and
$L_+\Lambda Q=-4Q$ follows by differentiating the scaled equation.

For the branch, set
\[
 \mathcal F(U,\omega,v)=\Delta^2U+\omega U+iv\cdot\nabla U-F(U).
\]
This map is $C^2$ from $H^4\times\mathbb R^{1+N}$ to $L^2$.
Let $g_0=iQ$, $g_j=\partial_jQ$, and augment the profile equation by
\[
 \mathcal F(U,\omega,v)=\sum_{j=0}^N b_jg_j,
 \qquad \langle U-Q,g_j\rangle_{\mathbb R}=0.
\]
Its derivative with respect to $(U,b)$ is invertible: the real
selfadjoint Hessian at $Q$ is Fredholm, its kernel is exactly the
span of the $g_j$, and it has a bounded complementary inverse from
$L^2$ to $H^4$. The implicit function theorem gives a $C^2$ augmented
branch. Phase and translation invariance imply
\[
 \langle\mathcal F(U,\omega,v),iU\rangle_{\mathbb R}=0,
 \qquad
 \langle\mathcal F(U,\omega,v),\partial_kU\rangle_{\mathbb R}=0.
\]
The pairing matrix between the fixed $g_j$ and these moving symmetry
vectors is close to the invertible Gram matrix at $Q$. Hence $b=0$,
and the branch consists of exact profiles.

The gauge determines profiles uniquely near $Q$ up to phase and
translation. Mass-preserving scaling maps parameters to
$(\kappa^4\omega,\kappa^3v)$; after returning to the gauge, uniqueness
shows that the branch is scaling-closed. Its reduced action is $C^2$
and has the usual charge derivatives. Differentiation at $(1,0)$ gives
$D_{v_j}U=-iL_-^{-1}\partial_jQ$ in the fixed gauge.

Finally, the Hessians have common domain $H^4$. The potentials converge
in $L^\infty$, and $iv\cdot\nabla$ is small from $H^4$ to $L^2$.
Norm-resolvent convergence therefore preserves the rank $N+1$ of a
small isolated spectral cluster about zero. The $N+1$ independent
vectors $iU,\partial_jU$ fill that cluster. They belong to the operator
domain by the profile equation and elliptic bootstrapping. Thus there
are no further zero or near-zero eigenvalues, and the inverse on the
kernel complement, extended by zero, is norm-continuous on $L^2$.
This verifies the complementary-gap assumption rather than adding it.
\end{proof}

\subsection{The linear group bound from one negative direction}

\begin{lemma}\label{lem:r12-group-bound}
For each fixed sufficiently small $v\ne0$,
\[
 n_-(\mathcal L_v)=1,
 \qquad
 \|e^{t\mathcal A_v}\|_{\mathcal B(H^2)}\le C_v e^{\lambda_vt}
 \quad(t\ge0).
\]
On an invariant complement to the two real eigenvectors the group
has at most linear growth.
\end{lemma}
\begin{proof}
At $v=0$, sharp critical optimality gives nonnegativity of $L_+$ on
$Q^\perp$ and
\[
 \langle Q,L_+Q\rangle=-\frac8N\|Q\|_{2+8/N}^{2+8/N}<0.
\]
Since $L_-\ge0$, the full Hessian has exactly one negative direction.
For small $v$, the free operator
\[
 H_v=\Delta^2+1+iv\cdot\nabla
\]
has symbol $|\xi|^4+1-v\cdot\xi$ uniformly comparable to
$1+|\xi|^4$. The potential $DF(U_v)$ is a decaying bounded real
symmetric matrix and converges in $L^\infty$ to $DF(Q)$. The drift
is small from $H^4$ to $L^2$. This gives norm-resolvent convergence
with a common lower bound and a uniformly positive essential spectrum.
The isolated zero cluster contains exactly the prescribed $N+1$
symmetry modes. Hence $n_-(\mathcal L_v)=1$ for small $v$.
This concrete perturbation argument excludes negative spectrum arriving
from infinity, as well as from zero.

Fix $v$ and omit its subscript. The profile equation and $F(U)\in H^2$
give $U\in H^6$. The matrix multiplier $DF(U)$ is bounded on $H^2$.
The free operator $-iH_v$ generates an isometric group on $H^2$;
its bounded real-linear perturbation $iDF(U)$ therefore generates a
$C^0$ group there. The $L^2$ eigenvectors of $\mathcal A$ lie in
$H^4$, and their eigenvalue equations improve them to $H^6$.

Let $\mathcal A Y_\pm=\pm\lambda Y_\pm$, $\lambda>0$, and write
$B(f,g)=\langle\mathcal Lf,g\rangle$ for the real form. The identity
\[
 B(\mathcal Af,g)+B(f,\mathcal Ag)=0
\]
gives $B(Y_\pm,Y_\pm)=0$. Furthermore
$b=B(Y_+,Y_-)\ne0$: otherwise their span would be a two-dimensional
totally isotropic space modulo $\ker\mathcal L$, impossible for a
form of negative index one. Their span is disjoint from the kernel,
as is seen by applying $\mathcal A$ and $\mathcal A^2$ to any
proposed relation. Hence
\[
 W=\operatorname{span}_{\mathbb R}\{Y_+,Y_-\},\qquad
 Z=W^{\perp_B},\qquad X=W\oplus Z
\]
is an invariant decomposition with bounded projection
\[
 \Pi_Wh=\frac{B(h,Y_-)}bY_++\frac{B(h,Y_+)}bY_-.
\]
The form is nonnegative on $Z$, because its sole negative direction
is already in $W$. Its null space in $Z$ is exactly
$K=\ker\mathcal L$: if $B(z,z)=0$, nonnegativity implies
$B(z,Z)=0$, while the definition gives $B(z,W)=0$.

Let $\Pi_K$ be the $L^2$ orthogonal projection onto $K$. Compactness
of the decaying potential relative to the positive free form gives
\begin{equation}\label{eq:r12-center-coercivity}
 B(z,z)\ge c_v\|(I-\Pi_K)z\|_{H^2}^2,
 \qquad z\in Z.
\end{equation}
Indeed, a violating sequence in $Z\cap K^\perp$, normalized in $H^2$,
has a weak limit. Weak lower semicontinuity of the free form and
compactness of the potential force its form value to be nonpositive.
The preceding null-space statement makes that limit zero. Compactness
then makes the potential term vanish, contradicting coercivity of the
free form and the normalization.

The form $B(e^{t\mathcal A}z,e^{t\mathcal A}z)$ is conserved, first
on the generator domain and then by density. It controls the component
orthogonal to $K$ through \eqref{eq:r12-center-coercivity}. For an
$L^2$ orthonormal basis $g_j$ of $K$, write
$h=z+k$ with $k=\Pi_Kh$. In the $H^{-2}$--$H^2$ pairing,
\[
 \frac d{dt}\langle h,g_j\rangle
 =-\langle\mathcal Lz,Jg_j\rangle,
 \qquad
 \left|\frac d{dt}\langle h,g_j\rangle\right|
       \le C_v\|z\|_{H^2}.
\]
This proves
\[
 \|e^{t\mathcal A}|_Z\|_{\mathcal B(H^2)}\le C_v(1+|t|).
\]
The group on $W$ is exactly $e^{\pm\lambda t}$. Absorbing $1+t$
into a $\lambda$-dependent multiple of $e^{\lambda t}$ proves the
lemma.
\end{proof}

\subsection{An ancient nonlinear orbit and exact conservation}

\begin{proof}[Proof of Theorem~\ref{thm:r12-nonlinear-instability}]
Lemma~\ref{lem:r12-branch} verifies the hypotheses of the calculation in Appendix~\ref{sec:r11-spectral} and hence supplies $\lambda_v$ and
$Y_{+,v}$. Fix $v$ and suppress it. Normalize $\|Y_+\|_X=1$. On continuous
functions $h:(-\infty,0]\to X$ use the norm
\[
 \|h\|_* =\sup_{t\le0}e^{-\lambda t}\|h(t)\|_X.
\]
For small $\varepsilon>0$, consider on the ball $\|h\|_*\le2\varepsilon$
the map
\begin{equation}\label{eq:r12-ancient-fixedpoint}
 \Phi(h)(t)=\varepsilon e^{\lambda t}Y_+
       +\int_{-\infty}^t e^{(t-s)\mathcal A}R(h(s))\,ds.
\end{equation}
The integral converges in $X$, since Lemma~\ref{lem:r12-group-bound}
and \eqref{eq:r12-remainder} give
\begin{equation}\label{eq:r12-ancient-error}
 \|\Phi(h)(t)-\varepsilon e^{\lambda t}Y_+\|_X
 \le \frac{C_v C_R2^{1+\beta}}{\beta\lambda}
             \varepsilon^{1+\beta}e^{(1+\beta)\lambda t}.
\end{equation}
By \eqref{eq:r12-remainder-lipschitz}, its weighted Lipschitz constant
is at most $C'_v\varepsilon^\beta/(\beta\lambda)$.
Choose $\varepsilon$ small enough for invariance of the ball and a
contraction constant below $1/2$. The fixed point is a nonlinear mild
solution on every finite subinterval: split the integral in
\eqref{eq:r12-ancient-fixedpoint} at a finite initial time and use
the group law. It satisfies \eqref{eq:r12-ancient-asymptotic} and
converges strongly to $U$ backward in time. Conservation and continuity
of $M,P,E$ therefore imply
\begin{equation}\label{eq:r12-exact-invariants}
 (M,P,E)(U+h(t))=(M,P,E)(U),\qquad t\le0.
\end{equation}

It remains to establish departure from the entire symmetry orbit.
The map $(\theta,y)\mapsto e^{i\theta}U(\cdot-y)$ is $C^2$ into
$H^2$, with injective differential and tangent $K$. It has no
nontrivial stabilizer: a nonzero phase-periodic translation contradicts
$U\in L^2$. Nor can translations tending to infinity approach $U$ in
$L^2$. Thus every orbit point sufficiently near $U$ belongs to this
local chart. Put
\[
 \eta_v=\operatorname{dist}_{H^2}(Y_+,K)>0.
\]
It is positive because $\mathcal A$ annihilates $K$ and has nonzero
eigenvalue on $Y_+$. Taylor expansion of the chart gives, for small
$a>0$,
\begin{equation}\label{eq:r12-orbit-transverse}
 \operatorname{dist}_{H^2}(U+aY_+,\mathcal O_v)
                          \ge\frac{\eta_v}{2}a.
\end{equation}
Indeed, any competing orbit point with comparable distance lies in the
chart with parameters $O(a)$; projecting its expansion off $K$ proves
the inequality. No existence of a global nearest orbit point is needed.

Shrink $\varepsilon_v$ further so that
\eqref{eq:r12-orbit-transverse} holds at $a=\varepsilon_v$, and the
right side of \eqref{eq:r12-ancient-error} at $t=0$ is at most
$\eta_v\varepsilon_v/4$. Then
\[
 \operatorname{dist}_{H^2}(U+h(0),\mathcal O_v)
                         \ge\eta_v\varepsilon_v/4.
\]
Set $\rho_v=\eta_v\varepsilon_v/8$. For $0<\delta<\varepsilon_v$
let $t_\delta=\lambda^{-1}\log(\delta/\varepsilon_v)$ and
$u_{\delta,v}=U+h(t_\delta)$. Its expansion and exact invariant values
follow from \eqref{eq:r12-ancient-error} and
\eqref{eq:r12-exact-invariants}. Its solution in the moving frame is
$U+h(t+t_\delta)$. The first time $\tau_{\delta,v}$ at which the
orbital distance equals $\rho_v$ satisfies
\[
 \lambda^{-1}\log\frac{\rho_v}{2\delta}
 \le\tau_{\delta,v}\le
 \lambda^{-1}\log\frac{\varepsilon_v}{\delta},
\]
because $\|h(t+t_\delta)\|_X\le2\delta e^{\lambda t}$ and the
distance at $t=-t_\delta$ exceeds $\rho_v$.
This proves \eqref{eq:r12-exit-time}. Formula \eqref{eq:r12-rate}
is the asymptotic in Theorem~\ref{thm:r11-spectral-pair} with $m=4$.
\end{proof}

\begin{remark}[Velocity dependence and the remaining assumption]
The construction gives a fixed positive escape radius for each fixed
nonzero small velocity. It does not bound that radius uniformly as
$v\to0$. The eigenvectors approach the neutral sector and the spectral
projection constants may diverge. The rigorously implied rate statement
along a unit direction $e$ is the iterated limit
\[
 \lim_{r\downarrow0}r\left[
   \lim_{\delta\downarrow0}
      \frac{\tau_{\delta,re}}{\log(1/\delta)}\right]
       =\sqrt{\frac{\mathcal D_Q}{2e^TTe}}.
\]
Critical scaling transfers the theorem to every frequency $\omega>0$,
with normalized velocity $v\omega^{-3/4}$ and the exact rate covariance
$\lambda(\omega,v)=\omega\lambda(1,v\omega^{-3/4})$.
The nonlinear instability mechanism is a classical Hamiltonian one;
the present connection uses the previously computed critical
small-velocity eigenvalue. Real-profile nondegeneracy is still an
explicit hypothesis. The theorem does not assert finite-time blowup or
classify the subsequent evolution after departure.
\end{remark}

\section{Conserved internal correlations and the critical mass}
\label{sec:r12-internal-dynamics}

The quantitative polarization theorem has a consequence for every time
of an actual solution. The internal correlation operator is conserved,
and a field with more than one occupied spatial mode cannot approach
the scalar sharp interpolation constant arbitrarily closely. This
raises the mass range in which conservation of energy controls the
whole evolution.

\subsection{The common scalar flow}
Let $\mathcal H$ be a nonzero separable complex Hilbert space and let
$a$ be a positive elliptic, real-even homogeneous symbol of order
$m>1$. Put $s=m/2$ and
\[
 \rho_U=\|U\|_{\mathcal H}^2,\qquad
 M(U)=\int\rho_U,\qquad
 A(U)=\int a(\xi)\|\widehat U(\xi)\|_{\mathcal H}^2\,d\xi.
\]
As in the quantitative polarization theorem, take either
$\mathcal F(U)=\int\rho_U^\kappa$, $\kappa=1+m/N$, or the positive
definite critical Hartree interaction
$\mathcal F(U)=\mathcal D(\rho_U,\rho_U)$, $\kappa=2$.
Write $r=\kappa-1$. We allow a common real nonnegative multiplier
$b$ of order at most $m$, and write
\begin{equation}\label{eq:r12-internal-flow}
 iU_t=[a(D)+b(D)]U-V_{\rho_U}U,\qquad
 E_b(U)=\frac{A(U)+B_b(U)}2-\frac{\mathcal F(U)}{2\kappa}.
\end{equation}
No homogeneity or dilation inequality on $b$ is used in this section.

For the general flow statements we assume a conservative local
$H^s(\mathbb R^N;\mathcal H)$ theory with the continuation property:
every finite endpoint of the maximal interval entails an unbounded
$H^s$ norm. A concrete quartic class satisfying this assumption is
proved below. The sharp scalar mass $M_c$ is fixed by
\[
 \mathcal F(u)\le\kappa(M(u)/M_c)^r A(u).
\]

Define the positive trace-class operator on $\mathcal H$
\begin{equation}\label{eq:r12-Gamma}
 \Gamma(U)=\int |U(x)\rangle\langle U(x)|\,dx,
 \qquad \mu(U)=\|\Gamma(U)\|_{\mathrm{op}},
 \qquad \chi(U)=1-\frac{\mu(U)}{M(U)}.
\end{equation}
For a nonzero field, $\operatorname{Tr}\Gamma=M$ and $0\le\chi<1$.
The largest eigenvalue exists even if the internal space is
infinite-dimensional, because $\Gamma$ is compact.

\begin{lemma}[Conservation of all internal correlations]
\label{lem:r12-Gamma-conservation}
Every energy-space solution of \eqref{eq:r12-internal-flow} satisfies
$\Gamma(U(t))=\Gamma(U(0))$. In particular, $\mu$ and $\chi$ are
constant throughout its maximal interval.
\end{lemma}
\begin{proof}
For fixed internal vectors $h,k$, the scalar components
$u_h=\langle h,U\rangle$ and $u_k=\langle k,U\rangle$ solve the
same equation $iu_t=H(t)u$, where
$H(t)=a(D)+b(D)-V_{\rho_U}$ is a self-adjoint spatial form.
The Hilbert-triple product rule for
$C H^s\cap C^1 H^{-s}$ functions gives
\[
 \frac{d}{dt}\int u_h\overline{u_k}
 =-i\langle H(t)u_h,u_k\rangle
       +i\langle u_h,H(t)u_k\rangle=0.
\]
The local nonlinear form is continuous by H\"older at exponent
$2\kappa$; the Hartree form is continuous by HLS. Equality for a
countable dense set of $h,k$ proves equality of the operators.
The trace-norm estimate
\[
 \|\Gamma(U)-\Gamma(W)\|_1
 \le(\|U\|_2+\|W\|_2)\|U-W\|_2
\]
also shows that this identity passes to energy-space limits.
\end{proof}

The conserved eigenvalues measure the masses in mutually orthogonal
spatial modes after an internal unitary change of basis. Merely writing
one spatial profile in many components still gives a rank-one operator
and $\mu=M$.

\subsection{A strict improvement retaining the scalar sharp mass}
\begin{theorem}[Internal mass distribution and energy control]
\label{thm:r12-mixed-threshold}
There is a constant $c>0$, depending only on the scalar dispersion
and interaction and independent of the number of components, such that
every nonzero $U\in H^s(\mathcal H)$ satisfies
\begin{equation}\label{eq:r12-mixed-GN}
 \mathcal F(U)
 \le \kappa\left(\frac{M(U)}{M_c}\right)^r
        [1-c\chi(U)]A(U).
\end{equation}
One may take $c=\min\{C^{-1},1/2\}$, where $C$ is the
dimension-independent quantitative polarization constant.

Consequently, if the initial state satisfies
\begin{equation}\label{eq:r12-mixed-positive}
 \eta:=1-\left(\frac{M}{M_c}\right)^r(1-c\chi)>0,
\end{equation}
its conservative energy-space evolution is global in both time
directions, and
\begin{equation}\label{eq:r12-mixed-global-bound}
 \sup_{t\in\mathbb R}\{A(U(t))+B_b(U(t))\}
 \le\frac{2E_b(U(0))}{\eta}.
\end{equation}
Thus every fixed $\chi>0$ permits the mass interval
\[
 M<M_c(1-c\chi)^{-1/r},
\]
which extends strictly above the scalar critical mass.
\end{theorem}
\begin{proof}
Let $\mathfrak d$ and $\mathfrak p$ denote the normalized
interpolation deficit and polarization distance:
\[
 \mathfrak d(U)=1-
 \frac{\mathcal F(U)}{\kappa(M(U)/M_c)^r A(U)},
 \qquad
 \mathfrak p(U)=\inf_{\|\zeta\|=1,q\in H^s}
 \left\{\frac{M(U-\zeta q)}{M(U)}
              +\frac{A(U-\zeta q)}{A(U)}\right\}.
\]
For fixed $\zeta$, orthogonal projection gives
\[
 \inf_q M(U-\zeta q)=M-\langle\zeta,\Gamma(U)\zeta\rangle.
\]
Taking the infimum over unit $\zeta$ and discarding the nonnegative
kinetic term yields $\mathfrak p(U)\ge\chi(U)$.
The quantitative polarization theorem gives
$\mathfrak p(U)\le C\mathfrak d(U)$, proving
\eqref{eq:r12-mixed-GN}.

For clarity, if
$\Gamma_A=\int a(\xi)|\widehat U(\xi)\rangle
\langle\widehat U(\xi)|\,d\xi$, the exact distance formula is
\[
 \mathfrak p(U)=2-\lambda_{\max}
       \left(\frac{\Gamma(U)}M+\frac{\Gamma_A}A\right).
\]
The same scalar projection minimizes both terms in their combined
Hilbert norm. Since $\Gamma_A/A$ is positive with trace one, this
also gives $\mathfrak p\ge\chi$. No conservation of $\Gamma_A$ is
asserted.

By Lemma~\ref{lem:r12-Gamma-conservation}, $M,\chi$, and hence
$\eta$ are constant along the flow. Equation
\eqref{eq:r12-mixed-GN} gives
\[
 E_b(U(t))\ge\frac\eta2A(U(t))+\frac12B_b(U(t)).
\]
As $0<\eta\le1$, conservation yields
\eqref{eq:r12-mixed-global-bound}. Ellipticity, the conserved mass,
and the continuation property finish the proof.
\end{proof}

\begin{corollary}[Scalar reduction of critical-mass unbounded evolution]
\label{cor:r12-minimal-scalar}
Suppose $M(U)=M_c$ and a conservative orbit has unbounded kinetic
energy along a sequence of times, including possible approach to a
finite maximal endpoint. Then $\Gamma(U(0))$ has rank one, and there
is one fixed unit vector $\zeta\in\mathcal H$ such that
\[
 U(t,x)=\zeta u(t,x)
\]
throughout the maximal interval. Here $u$ solves the corresponding
scalar equation with mass $M_c$.

More generally, any energy-space-unbounded orbit of finite conserved
mass satisfies
\begin{equation}\label{eq:r12-nearcritical-mixing}
 \chi\le c^{-1}\left[1-(M_c/M)^r\right].
\end{equation}
\end{corollary}
\begin{proof}
If $M=M_c$ and $\chi>0$, Theorem~\ref{thm:r12-mixed-threshold}
applies with $\eta=c\chi$, contradicting unboundedness. Thus
$\mu=M$, so positivity and trace $M$ force
$\Gamma=M|\zeta\rangle\langle\zeta|$. For every
$h\perp\zeta$,
\[
 \int|\langle h,U(t,x)\rangle|^2\,dx
     =\langle h,\Gamma h\rangle=0.
\]
A countable orthonormal basis in $\zeta^\perp$ shows that
$U(t,x)$ lies in the same complex line almost everywhere.
Substitution gives the scalar equation. Finally an unbounded orbit
cannot satisfy \eqref{eq:r12-mixed-positive}; rearranging its
negation proves \eqref{eq:r12-nearcritical-mixing}.
\end{proof}

The scalar reduction holds for the initial datum and the entire
orbit, not merely for a rescaled concentration limit. For
$M=M_c(1+\varepsilon)$, \eqref{eq:r12-nearcritical-mixing} gives
$\chi\le(r/c)\varepsilon+O(\varepsilon^2)$.

\subsection{A concrete quartic flow}
\begin{proposition}[Unconditional realization in dimensions one to three]
\label{prop:r12-quartic-vector-flow}
Let $N=1,2,3$, let $a$ be any real-even positive elliptic homogeneous
quartic symbol, and let $b$ be real with
$0\le b(\xi)\le C(1+|\xi|^4)$. Then
\[
 iU_t=[a(D)+b(D)]U-\|U\|_{\mathcal H}^{8/N}U,
 \qquad U(0)\in H^2(\mathbb R^N;\mathcal H),
\]
has a conservative local $H^2$ flow with the continuation property.
Therefore Theorem~\ref{thm:r12-mixed-threshold} and
Corollary~\ref{cor:r12-minimal-scalar} hold without additional flow
assumptions for these equations, including countably many components.
\end{proposition}
\begin{proof}
The Fourier group is unitary on $H^2$. In these dimensions,
$H^2$ embeds into $L^\infty$, and the radial Hilbert-space map
$z\mapsto\|z\|^{8/N}z$ is locally Lipschitz on $H^2(\mathcal H)$.
Indeed its first three derivatives are bounded on bounded balls
because $8/N>2$; the chain rule, the $H^2$ product estimates, and
$H^1\hookrightarrow L^4$ give the required estimate, with constants
independent of internal dimension. Duhamel contraction gives the
local flow and its continuation property. Symmetric Fourier
Galerkin regularization gives mass and energy conservation; the
regularized solutions converge in $C H^2$ on every local existence
interval, and both conserved functionals are continuous there.
Lemma~\ref{lem:r12-Gamma-conservation} gives conservation of the
full internal operator.
\end{proof}

\paragraph{Relation to kinetic-energy inequalities.}
Conservation of all component overlaps and the use of the largest
occupation are established tools in nonlinear Schr\"odinger systems;
see Hong--Kwon--Yoon, arXiv:1712.06811, Section~1 and
Proposition~7.3. Their Theorem~2.3 also includes fractional
kinetic-energy inequalities. The specific connection used here is
the dimension-independent correction $1-c(1-\mu/M)$ to the
\emph{scalar sharp} critical-mass inequality. This gives a strict
extension above $M_c$ for every genuinely nonpolarized internal
state and the exact initial-data reduction at mass $M_c$.
The constant $c$ is not claimed to be optimal, and equality in
\eqref{eq:r12-mixed-positive} is not covered.


\section{Relation to prior work and the remaining evolution problem}
Boulenger--Lenzmann~\cite{BL}, Theorem 3, prove finite- or infinite-time
blowup for radial negative-energy data in the pure critical quartic
equation, in dimensions at least two. Bonheure--Cast\'eras--Gou--Jeanjean
\cite{BCGJ}, Theorem 1.2, deduce the corresponding instability of radial
ground states, without real-sector nondegeneracy. The recent
Dinh--Muse manuscript~\cite{DM}, Theorem 1.1, states finite-time blowup
for radial negative-energy data, while its Remark 1.4 leaves the
nonradial problem open. Theorem~\ref{thm:r12-amplification} treats
unrestricted spatial geometry but adds zero momentum and small mass
excess, and concludes quantitative amplification rather than blowup.
It is not a strengthening of the radial finite-time conclusion.

Instability obtained with radial perturbations already implies
instability in the full perturbation space of that same radial wave.
Moreover, Lenzmann--Sok~\cite{LS}, Theorem 2, prove radial symmetry
of all sharp optimizers when the Lebesgue exponent is even. In the
quartic critical problem this includes $N=1,2,4$; the radial instability
results above therefore already cover all ground states for $N=2,4$.
The distinct assertion here is the quantitative constraint on arbitrary
nonradial near-threshold zero-momentum trajectories. The targeted
comparison found no matching $d^{5/3}$ estimate, but does not establish
exhaustive priority. Gou's cylindrical result~\cite{Gou} imposes both
cylindrical symmetry and an axial moment condition and does not give
that unrestricted statement.

Existence of traveling ground waves for general dispersion is studied
by Bugiera--Lenzmann--Schikorra--Sok~\cite{BLSS}. We include the
small-velocity ground-wave compactness proof needed for constrained
descent. Hamiltonian unstable subspaces and group bounds are established
tools; Lin--Zeng~\cite{LZ} provide a general trichotomy theory.
Our nonlinear construction uses these ideas in a concrete form to turn
the computed critical small-velocity eigenvalue into an ancient orbit
and an exit time. Real $L_+$ nondegeneracy remains an assumption for
that result, and is not supplied by the universally valid phase gap.

Conservation of mutual component inner products is established in
unitary nonlinear Schr\"odinger systems, including
Hong--Kwon--Yoon~\cite{HKY}, Proposition 7.3. Their generalized
kinetic-energy inequalities also include fractional orders and the
relevant endpoint. We do not claim those tools as new. The connection
used here retains the sharp scalar $M_c$, introduces a quantitative
loss $1-c\chi$ through polarization, and then uses the conserved
$\chi$ to control every time of the evolution. The resulting enlarged
boundedness interval is sufficient; it is not claimed to be the exact
mixed-state critical mass.

\paragraph{What is still missing.}
A lower bound on the largest kinetic norm does not exclude a global
bounded solution that reaches that large norm. Excluding such a
trajectory, and then distinguishing finite-time concentration from
infinite-time growth, remain separate unresolved steps. Zero momentum
does not set the fourth-order mass-transport velocity to zero: that
velocity involves a cubic Fourier moment. A sharper long-time link
between conserved momentum, core motion, and changes of scale is the
next central issue. We have not proved that link here.

\appendix
\section{Sharp density geometry and quantitative polarization}
\label{sec:r12-polarization-foundation}
This appendix reproduces the analytic inputs used by the new evolution
arguments. They are proved here, including the infinite-dimensional
compactness step. No radiality, positivity of the kinetic semigroup,
or real-sector nondegeneracy is used.

Let $\mathcal H$ be any nonzero separable complex Hilbert space.
Finite-dimensional $\mathbb C^d$ and $\ell^2$ are both allowed.
Fix $N\ge1$, $m>1$, $s=m/2$, and a real-even, positive elliptic,
$m$-homogeneous symbol $a$, of class $C^{N+2}$ away from zero.
For $U\in H^s(\mathbb R^N;\mathcal H)$ put
\[
 \rho_U=\|U\|_{\mathcal H}^2,\quad
 M(U)=\int\rho_U,\quad
 A(U)=\int a(\xi)\|\widehat U(\xi)\|_{\mathcal H}^2\,d\xi.
\]
We use either of the following interactions:
\begin{align*}
 &\kappa=1+m/N,
 &&\mathcal F(U)=\int\rho_U^\kappa,
 &&V_\rho=\rho^{\kappa-1};\\
 &\kappa=2,
 &&\mathcal F(U)=\mathcal D(\rho_U,\rho_U),
 &&V_\rho=K*\rho.
\end{align*}
In the second case $1<m<N$, the real even kernel $K$ is homogeneous
of degree $-m$, its interaction form $\mathcal D$ is positive definite,
and
\[
 |K(x)|+|x||\nabla K(x)|\le C|x|^{-m}.
\]
As before, positivity of the Fourier transform of $K$ almost everywhere
suffices for positive definiteness. Pointwise positivity is not needed.
The scalar sharp inequality defines $M_c$ by
\begin{equation}\label{eq:r11-vectorGN}
 \mathcal F(u)\le\kappa
       \left(\frac{M(u)}{M_c}\right)^{\kappa-1}A(u).
\end{equation}

\begin{theorem}[Sharp interpolation and dimension-independent compactness]
\label{thm:r11-vectorcompact}
Inequality \eqref{eq:r11-vectorGN} has the same sharp constant for
$\mathcal H$-valued functions. Every optimizer has the form
\[
 U(x)=\zeta Q(x),\qquad \|\zeta\|_{\mathcal H}=1,
\]
where $Q$ is a real scalar optimizer, up to a constant phase.
If $M(U_n)\to M_c$, $A(U_n)\to1$, and
$\mathcal F(U_n)\to\kappa$, then the sequence is strongly compact
in $H^s$ after translations and internal unitaries. This remains
true when the internal dimensions vary: embed all spaces in
$\ell^2$ before applying the unitaries.
\end{theorem}

\subsection{Sharp interpolation and complete polarization}
The local interaction is the $\kappa$-th power of the $L^\kappa$
density norm. The Hartree interaction is the square of its positive
definite Hilbert density norm. Thus both have the form
$\mathcal F(U)=\|\rho_U\|_{\mathcal B}^\kappa$ for a strictly convex
real density space $\mathcal B$.
In an internal orthonormal basis $U=(u_j)$, the triangle inequality,
scalar sharp interpolation, and sequence H\"older give
\begin{align*}
 \mathcal F(U)^{1/\kappa}
 &\le\sum_j\mathcal F(u_j)^{1/\kappa}\\
 &\le\kappa^{1/\kappa}M_c^{-(\kappa-1)/\kappa}
   \sum_j A(u_j)^{1/\kappa}M(u_j)^{(\kappa-1)/\kappa}\\
 &\le\kappa^{1/\kappa}M_c^{-(\kappa-1)/\kappa}
              A(U)^{1/\kappa}M(U)^{(\kappa-1)/\kappa}.
\end{align*}
Finite partial sums converge in the density norm because the same
bound applied to their tails tends to zero. A scalar optimizer placed
in one component proves sharpness.

Split the components further into real scalar functions $f_j$.
Equality forces every active $f_j$ to be a scalar optimizer, all
$f_j^2$ to be proportional, and all kinetic-to-mass ratios to agree.
The same conclusions hold after any real orthogonal rotation of a
pair $f_i,f_j$. If $f_j/f_i$ took both signs of a nonzero constant,
the squared sum and difference could not be proportional; when the
constant has magnitude one their nonzero supports would be disjoint.
Thus each $f_j$ is one real constant multiple of a fixed scalar
function. This proves complete polarization, including for countably
many components.

\subsection{Compactness modulo internal unitaries}
Ordinary local compactness does not hold directly in an infinite
internal dimension: a fixed spatial shape can move through mutually
orthogonal internal vectors. The following argument removes precisely
that loss of compactness.

\begin{lemma}[Internal-unitary local compactness]
\label{lem:r11-internalcompact}
Let $X=H^s(\mathbb R^N;\mathbb C)$ with $s>0$. From every bounded
sequence in $X\widehat\otimes\ell^2$, one can extract a subsequence
and internal unitaries $W_n$ such that $W_nU_n$ converges weakly in
$X\widehat\otimes\ell^2$ and strongly in
$L^2_{\mathrm{loc}}(\mathbb R^N;\ell^2)$.
\end{lemma}
\begin{proof}
Fix an orthonormal basis $(\phi_k)$ of the scalar Hilbert space $X$.
Write its internal coefficient vectors as
$c_{nk}=\langle\phi_k,U_n\rangle_X\in\ell^2$.
Gram--Schmidt the first $n$ vectors and extend the resulting finite
orthonormal list to an internal orthonormal basis. A unitary $W_n$
can therefore be chosen so that
\[
 W_nc_{nk}\in\operatorname{span}\{e_1,\ldots,e_k\},\qquad k\le n.
\]
For each fixed $k$, these vectors have a convergent subsequence in a
fixed finite-dimensional space. Diagonal extraction gives strong
convergence of every coefficient vector. Weak compactness gives a
weak limit of $W_nU_n$, with exactly those coefficients.

For $\chi\in C_c^\infty$, the scalar multiplication map
$K_\chi:X\to L^2$, $f\mapsto\chi f$, is compact. If $P_L$ is the
scalar projection onto the first $L$ basis vectors, then
\[
 \|K_\chi(I-P_L)\|_{X\to L^2}\longrightarrow0.
\]
Tensoring with the identity on $\ell^2$ preserves this norm.
The transformed $P_LU_n$ converge strongly because there are only
$L$ coefficient vectors; after multiplication by $\chi$, the tails
are uniformly small. This proves local strong convergence.
\end{proof}

\begin{proof}[Compactness assertion of Theorem~\ref{thm:r11-vectorcompact}]
The localized Sobolev inequality has constants independent of internal
dimension: use the Hilbert-valued Bessel-potential representation and
its norm triangle inequality, then spatial localization and the sum of
component quadratic norms. Thus local mass vanishing for an
$H^s$-bounded sequence implies vanishing in $L^p(\mathcal H)$ for
$2<p<2N/(N-2s)$, or every finite $p>2$ if $2s\ge N$.
In the local case take $p=2\kappa$; in the Hartree case use
$p=4N/(2N-m)$ and Hardy--Littlewood--Sobolev. Hence
$\mathcal F(U_n)\to\kappa>0$ excludes local mass vanishing.

Translate a nonzero amount of local mass into a fixed ball and use
Lemma~\ref{lem:r11-internalcompact}. We obtain a nonzero weak limit
$U$, strong local $L^2$ convergence, and by interpolation strong local
convergence in the interaction exponent. Put $R_n=U_n-U$.
The quadratic quantities split. For local powers the Brezis--Lieb
argument for $\|U\|_{\mathcal H}^{2\kappa}$ gives
\[
 \mathcal F(U_n)=\mathcal F(U)+\mathcal F(R_n)+o(1).
\]
For Hartree, the term $\operatorname{Re}\langle U,R_n\rangle$
tends to zero in $L^{p/2}$ by local convergence and the small
$L^p$ tail of $U$. Also
$\mathcal D(\rho_U,\rho_{R_n})\to0$: approximate $U$ in
$L^p(\mathcal H)$ by a compactly supported smooth finite-dimensional
field. Its potential is bounded and decays at infinity, while $R_n$
tends locally to zero in $L^2$. HLS controls the approximation error.
This proves the same splitting for Hartree.

Let $M_1>0,M_2$ and $A_1,A_2$ be the limiting contributions.
Then $M_1+M_2=M_c$, $A_1+A_2=1$, and sharp interpolation yields
\[
 M_c^{\kappa-1}\le
 M_1^{\kappa-1}A_1+M_2^{\kappa-1}A_2.
\]
Both masses cannot be positive, since each would be strictly below
$M_c$. Thus $M_2=0$, and the inequality then forces $A_2=0$.
The convergence is strong in $H^s$. The limit is an optimizer,
already proved to be completely polarized.
\end{proof}

\subsection{Phase coercivity and the internal symmetry kernel}
For completeness, the scalar phase fact used below follows directly
from the density norm. Normalize a real optimizer by $M(Q)=M_c$,
$A(Q)=1$, and put $\omega_*=(\kappa-1)/M_c$. The form
\[
 L_Q=a(D)+\omega_*-V_{Q^2}
\]
is nonnegative: differentiate the sharp inequality at
$Q+i\sqrt t\,h$. If $h$ belongs to its kernel, write
$x=A(h)$ and $y=M(h)/M_c$. The supporting-functional inequality
for the density norm, sharp interpolation for $h$, and weighted
arithmetic--geometric mean give
\[
 x+(\kappa-1)y
 \le\kappa x^{1/\kappa}y^{(\kappa-1)/\kappa}
 \le x+(\kappa-1)y.
\]
Equality implies $x=y$ and $h^2=yQ^2$ by strict convexity.
The kernel of a nonnegative form is linear; apply the same conclusion
to $h+Q$ to get $h=cQ$. The potential is compact as a form
$H^s\to H^{-s}$, by local Rellich and H\"older in the local case,
or the compact-potential approximation and HLS in the Hartree case.
Consequently
\begin{equation}\label{eq:r11-vectorcoercivity}
 \langle h,L_Qh\rangle\ge c_Q\|h\|_{H^s}^2,
 \qquad \langle h,Q\rangle_{L^2}=0.
\end{equation}
This holds for complex $h$ by real--imaginary splitting, and it can
be summed over arbitrarily many components with the same constant.

At $Qe_1$, write a perturbation as $(h+ik)e_1+Z$ with $Z\perp e_1$.
The Hessian of $E+\omega_*M/2$ is
\[
 \langle h,L_+h\rangle+\langle k,L_Qk\rangle
 +\sum_{j\ge2}\bigl(
  \langle\operatorname{Re}Z_j,L_Q\operatorname{Re}Z_j\rangle
 +\langle\operatorname{Im}Z_j,L_Q\operatorname{Im}Z_j\rangle\bigr).
\]
Its kernel in the phase and internal transverse sector is exactly
\[
 \{Q\eta:\ \eta\in\mathcal H,
                 \operatorname{Re}\langle\eta,e_1\rangle=0\}.
\]
For $\mathcal H=\mathbb C^d$, this is the $(2d-1)$-dimensional real
tangent space to the internal unitary orbit. Any additional full
Hessian kernel lies in the scalar real sector $L_+$; no assumption
on that sector is made.

\subsection{Quantitative polarization without real-sector nondegeneracy}
The preceding compactness and phase coercivity also give a quantitative
statement for arbitrary fields, not just stationary or traveling ones.
Define the scale-invariant interpolation deficit
\[
 \mathfrak d(U)=1-
 \frac{\mathcal F(U)}{\kappa(M(U)/M_c)^{\kappa-1}A(U)}
\]
and the squared distance to a completely polarized field by
\[
 \mathfrak p(U)=\inf_{\substack{\|\zeta\|_{\mathcal H}=1\\q\in H^s}}
 \left\{\frac{M(U-\zeta q)}{M(U)}
             +\frac{A(U-\zeta q)}{A(U)}\right\}.
\]
Both are defined for every nonzero $U\in H^s(\mathcal H)$.

\begin{theorem}[Dimension-independent quantitative polarization]
\label{thm:r11-quantitativepolarization}
There is $C<\infty$, depending only on the fixed scalar dispersion
and interaction, such that
\begin{equation}\label{eq:r11-polarizationdeficit}
 \mathfrak p(U)\le C\mathfrak d(U)
\end{equation}
for every nonzero Hilbert-valued field $U$ and every internal dimension.
No nondegeneracy or uniqueness of the scalar real-sector ground state
is required.
\end{theorem}
\begin{proof}
Amplitude and spatial scaling preserve both quantities, so impose
$M(U)=M_c$, $A(U)=1$. Write
$\|U\|_X^2=M(U)/M_c+A(U)$.
Since $\mathfrak p(U)\le2$, failure of a uniform constant would
produce $U_n$ with $\mathfrak d(U_n)\to0$ and
$\mathfrak d(U_n)/\mathfrak p(U_n)\to0$.
After the symmetries of Theorem~\ref{thm:r11-vectorcompact},
$U_n\to Qe_1$ strongly in $H^s$.
Align $\int QU_n$ with the positive real first internal direction
and write $U_n=(u_n,Z_n)$. Each component of $Z_n$ is complex
$L^2$-orthogonal to $Q$, and $u_n\to Q$, $Z_n\to0$ in $H^s$.

The interaction expansion in these transverse directions is
\begin{equation}\label{eq:r11-transverseexpansion}
 \mathcal F(u_n,Z_n)=\mathcal F(u_n)
 +\kappa\langle V_{|u_n|^2},\|Z_n\|^2\rangle
 +o(\|Z_n\|_X^2).
\end{equation}
For local powers, set $a=|u_n|^2$ and $b=\|Z_n\|^2$ pointwise.
When $1<\kappa\le2$, the nonnegative Taylor remainder of
$(a+b)^\kappa$ after its linear term in $b$ is at most $C_\kappa b^\kappa$.
When $\kappa>2$, it is at most
$C_\kappa(a^{\kappa-2}b^2+b^\kappa)$.
H\"older and the dimension-independent $L^{2\kappa}$ Sobolev bound
make these respectively $O(\|Z_n\|_X^{2\kappa})$ and
$O(\|Z_n\|_X^4+\|Z_n\|_X^{2\kappa})$, since $u_n$ stays bounded.
Both are $o(\|Z_n\|_X^2)$.
For Hartree the remainder is exactly
$\mathcal D(\|Z_n\|^2,\|Z_n\|^2)=O(\|Z_n\|_X^4)$.

Put $m_n=M(Z_n)$ and $a_n=A(Z_n)$. Scalar sharp interpolation gives
\[
 \mathcal F(u_n)\le
 \kappa(1-m_n/M_c)^{\kappa-1}(1-a_n).
\]
Combining its expansion with \eqref{eq:r11-transverseexpansion}
and potential form convergence yields
\begin{align*}
 \kappa-\mathcal F(U_n)
 &\ge\kappa\left[
 a_n+\frac{\kappa-1}{M_c}m_n
       -\langle V_{Q^2},\|Z_n\|^2\rangle\right]
       -o(\|Z_n\|_X^2)\\
 &=\kappa\sum_j\langle Z_{nj},L_QZ_{nj}\rangle
       -o(\|Z_n\|_X^2)\\
 &\ge c\|Z_n\|_X^2.
\end{align*}
The constant is independent of component number. Taking the candidate
$e_1u_n$ in the distance gives
$\mathfrak p(U_n)\le\|Z_n\|_X^2$.
As $\mathfrak d(U_n)=(\kappa-\mathcal F(U_n))/\kappa$,
this contradicts the assumed ratio and proves the theorem.
\end{proof}

The theorem controls internal polarization rather than the full
scalar shape. It therefore remains valid even when scalar optimizer
uniqueness or real-sector quantitative stability is unavailable.


\section{The spectral pair created by slow translation}
\label{sec:r11-spectral}

The velocity--momentum identities determine the inertia of the charge Hessian. We now identify the pair of spectral values that emerges
from the critical scaling mode, including its leading coefficient. The
result assumes real-profile nondegeneracy; this hypothesis is not supplied
by phase coercivity alone.

Let $a$ be positive, elliptic, real-even and homogeneous of degree $m>1$.
Consider
\[
 iu_t=a(D)u-|u|^{2m/N}u,
 \qquad
 a(D)U+\omega U+iv\cdot\nabla U-|U|^{2m/N}U=0.
\]
All Hilbert spaces in this section are real unless explicitly complexified.
Let $Q=U_{1,0}$ be a real sharp ground wave and set
\[
 \Lambda Q=\frac N2Q+x\cdot\nabla Q,\qquad
 L_-=a(D)+1-|Q|^{2m/N},\qquad
 L_+=a(D)+1-\left(1+\frac{2m}{N}\right)|Q|^{2m/N}.
\]
We impose the following hypotheses.
\begin{enumerate}
\item The functions $\partial_jQ$ and $\Lambda Q$ have the regularity and
spatial decay needed to differentiate translations and critical scaling
in $L^2$ and in the profile equation. In particular,
$L_+\Lambda Q=-mQ$ holds in the operator sense.
\item $L_-\geq0$, $\ker L_-=\mathbb RQ$, and $L_-^{-1}$ is bounded
on $Q^\perp$. Moreover
$\ker L_+=\operatorname{span}\{\partial_jQ:1\leq j\leq N\}$,
and $L_+$ has a bounded inverse on that kernel's orthogonal complement.
\item There is a $C^1$ local profile branch $U_{\omega,v}$, closed under
critical scaling up to phase and translation. Its reduced action
$d(\omega,v)=E(U_{\omega,v})+\omega M(U_{\omega,v})/2
-v\cdot P(U_{\omega,v})/2$ is $C^2$ and satisfies
$d_\omega=M/2$ and $d_v=-P/2$. Profile differentiation is valid, and in
a translation gauge
\[
 D_{v_j}U_{1,0}=-iL_-^{-1}\partial_jQ.
\]
\item On the real space $\mathcal H=L^2(\mathbb R^N;\mathbb C)$ the
real action Hessians $\mathcal L_v$ are selfadjoint on a common dense
domain. Their kernels are exactly
\[
 K_v=\operatorname{span}_{\mathbb R}
       \{iU_{1,v},\partial_1U_{1,v},\ldots,\partial_NU_{1,v}\}.
\]
These kernel bases vary continuously in $\mathcal H$. The inverse of
$\mathcal L_v$ on $K_v^\perp$, extended by zero on $K_v$, exists and is
norm-continuous as a bounded operator on $\mathcal H$ for small $v$.
\end{enumerate}
The last hypothesis explicitly includes a uniform complementary spectral
gap. Common-domain operator continuity, or suitable norm-resolvent
continuity with an isolated kernel of fixed dimension, provides standard
ways to verify it. It is not being inferred for every dispersion symbol.

Define the real matrix and vector
\[
 T_{ij}=\langle\partial_iQ,L_-^{-1}\partial_jQ\rangle,
 \qquad c_j=\langle\Lambda Q,L_-^{-1}\partial_jQ\rangle,
\]
and the scalar
\begin{equation}\label{eq:r11-spectral-D}
 \mathcal D_Q=
 \langle\Lambda Q,L_-^{-1}\Lambda Q\rangle-c^TT^{-1}c.
\end{equation}
The inverse of $L_-$ always acts on $Q^\perp$.

\begin{theorem}[Small-velocity spectral asymptotics]
\label{thm:r11-spectral-pair}
Under the four hypotheses above, $T$ is positive definite and
$\mathcal D_Q>0$. If $m\ne2$, then for every sufficiently small nonzero
$v$, the nonzero spectrum of
\[
 \mathcal A_v=J\mathcal L_v,\qquad J=-i,
\]
in a fixed neighborhood of zero consists of exactly one simple pair
$\lambda_\pm(v)$, satisfying
\begin{equation}\label{eq:r11-lambda-square}
 \lambda_\pm(v)^2
  =\frac{(m-2)v^TTv}{\mathcal D_Q}+o(|v|^2).
\end{equation}
For $m>2$ this pair is real and
\[
 \lambda_\pm(v)=
 \pm\sqrt{\frac{(m-2)v^TTv}{\mathcal D_Q}}+o(|v|).
\]
In particular, the wave is spectrally unstable. For $1<m<2$ the pair is
purely imaginary, with
\[
 \lambda_\pm(v)=
 \pm i\sqrt{\frac{(2-m)v^TTv}{\mathcal D_Q}}+o(|v|).
\]
The latter assertion concerns the pair near zero and does not, by itself,
exclude unstable spectrum elsewhere.
\end{theorem}

\begin{proof}
Phase coercivity gives $T>0$: a nonzero directional derivative of a
nonzero $L^2$ function cannot vanish. Put
\[
 \Lambda_\perp Q=
 \Lambda Q-\sum_j(T^{-1}c)_j\partial_jQ.
\]
Then
\[
 \mathcal D_Q=
 \langle\Lambda_\perp Q,L_-^{-1}\Lambda_\perp Q\rangle\geq0.
\]
Equality would imply
$(x-\zeta)\cdot\nabla Q+NQ/2=0$ for some $\zeta$.
Thus $Q$ would be fixed by the mass-preserving dilation group about
$\zeta$. A nonzero function homogeneous of degree $-N/2$ cannot belong
to $L^2$: its radial mass integral is proportional to $\int_0^\infty dr/r$.
Hence $\mathcal D_Q>0$. The same Schur-complement formula shows that
$\mathcal D_Q$ is unchanged if a linear combination of translations is
added to $\Lambda Q$, so it is independent of the spatial center.

Write $U_v=U_{1,v}$ and
$g_0=iU_v$, $g_j=\partial_jU_v$. Their span $K_v$ is isotropic:
\begin{equation}\label{eq:r11-isotropic-kernel}
 \langle g_i,Jg_j\rangle=0.
\end{equation}
For the phase--translation pairings this follows from integration of a
mass derivative. For two translations it follows from commutation of
translation charges, or directly from the Fourier representation of
$\langle\partial_iU,-i\partial_jU\rangle_{\mathbb R}$.

Let $\Pi_v$ project orthogonally onto $K_v$, let $\Pi_c=I-\Pi_v$, and
write $\mathcal L_c=\mathcal L_v|_{K_v^\perp}$ and
$J_c=\Pi_cJ\Pi_c$. The operator
\[
 R_v(\lambda)=(\mathcal L_c+\lambda J_c)^{-1}
\]
is extended by zero on $K_v$. A Neumann series and hypothesis 4 show
that this inverse is analytic for $|\lambda|<\rho$, with a common
$\rho>0$ for small $v$. The eigenvalue problem is considered on the
complexification of $\mathcal H$; real pairings below extend bilinearly.

The equation $J\mathcal L_v\psi=\lambda\psi$ is equivalent to
$(\mathcal L_v+\lambda J)\psi=0$. Decompose $\psi=k+w$ with
$k\in K_v$ and $w\in K_v^\perp$. For $\lambda\ne0$ the complementary
equation gives
\[
 w=-\lambda R_v(\lambda)Jk.
\]
By \eqref{eq:r11-isotropic-kernel}, the kernel equation is therefore
exactly $D(\lambda,v)z=0$, where
\begin{equation}\label{eq:r11-reduced-matrix}
 D_{ij}(\lambda,v)
  =\langle g_i,J R_v(\lambda)Jg_j\rangle,
 \qquad k=\sum_{j=0}^Nz_jg_j.
\end{equation}
The common factor suppressed in this equation is $-\lambda^2$.
For $0<|\lambda|<\rho$, analytic block elimination therefore preserves
the eigenvalues and their algebraic multiplicities. Selfadjointness of
$\mathcal L_c$ and skew-adjointness of $J$ imply
\begin{equation}\label{eq:r11-reduced-transpose}
 D(-\lambda,v)=D(\lambda,v)^T.
\end{equation}

Profile differentiation yields
\[
 \mathcal L_v\partial_\omega U=-U,
 \qquad
 \mathcal L_v\partial_{v_j}U=-i\partial_jU=Jg_j.
\]
Components of the parameter derivatives in $K_v$ do not contribute,
by \eqref{eq:r11-isotropic-kernel}. Consequently
\[
 D(0,v)=C d''(1,v)C,
 \qquad C=\operatorname{diag}(1,-I_N).
\]
The lower-right block is $D_{vv}(0,v)=-B(1,v)/2$, where
$B=D_vP$. The first derivative of the branch gives
\[
 P(U_v)=2Tv+o(|v|),\qquad B(1,v)=2T+o(1).
\]
Thus the lower-right block stays invertible for small $\lambda,v$.
Take its scalar Schur complement $f(\lambda,v)$. Identity
\eqref{eq:r11-reduced-transpose} makes $f$ even in $\lambda$.
It is consequently analytic in $z=\lambda^2$ in a common disk, with
coefficients continuous in $v$.

The exact scaling identity for the reduced action gives
\begin{align}
 f(0,v)
 &=\frac{P^TB^{-1}P-(m-1)v\cdot P}{2m^2}\notag\\
 &=\frac{2-m}{m^2}v^TTv+o(|v|^2).
 \label{eq:r11-f-zero}
\end{align}
For completeness, the identity follows by differentiating the scaling
laws for the charges. With $H=d''$ and
$D_*=(m\omega,(m-1)v)^T$, they give
$HD_*=(0,-P/2)^T$ and $H_{vv}=-B/2$. Eliminating the velocity block
shows that its scalar Schur complement is
$[P^TB^{-1}P-(m-1)v\cdot P]/(2m^2\omega^2)$.
This uses no differentiated second-order remainder for the profile.

It remains to compute the coefficient of $\lambda^2$ at $v=0$.
Let $R_0$ be the $L^2$-orthogonal projection of $\Lambda Q/m$ off the
translation kernel, and put $Z_j=L_-^{-1}\partial_jQ$.
At $v=0$ the two real blocks give
\[
 \mathcal L_0=\operatorname{diag}(L_+,L_-),
 \qquad J(h,b)=(b,-h),
\]
\[
 \mathcal L_c^{-1}Jg_0=(-R_0,0),
 \qquad \mathcal L_c^{-1}Jg_j=(0,-Z_j).
\]
Expanding the exact resolvent in \eqref{eq:r11-reduced-matrix}, and
writing $b_j=\langle R_0,Z_j\rangle$, gives
\begin{align*}
 D_{00}(\lambda,0)
  &=\lambda^2\langle R_0,L_-^{-1}R_0\rangle+O(\lambda^4),\\
 D_{0v}(\lambda,0)&=\lambda b^T+O(\lambda^2),\\
 D_{v0}(\lambda,0)&=-\lambda b+O(\lambda^2),\\
 D_{vv}(\lambda,0)&=-T+O(\lambda).
\end{align*}
The scalar Schur complement is even. Its coefficient is therefore
\[
 \partial_z f(0,0)=
 \langle R_0,L_-^{-1}R_0\rangle-b^TT^{-1}b
 =\frac{\mathcal D_Q}{m^2}>0.
\]
The final equality uses invariance of the Schur expression under
addition of translations to the scaling derivative. Analyticity with
a common radius and continuity in $v$ now give
\[
 f(\lambda,v)=\frac{2-m}{m^2}v^TTv+o(|v|^2)
 +\lambda^2\left(\frac{\mathcal D_Q}{m^2}+o(1)\right)
 +O(\lambda^4),
\]
where the last remainder is uniform for small $v$.

The analytic scalar equation, regarded as an equation in $z=\lambda^2$,
has one simple real root near zero, by
$\partial_zf(0,0)>0$ and uniform analytic inversion. That root obeys
\eqref{eq:r11-lambda-square}. Its sign is the sign of $m-2$, since
$T>0$, and it is nonzero for nonzero sufficiently small $v$.
Its two square roots are simple because
$\partial_\lambda f=2\lambda\partial_zf\ne0$ there.
The invertible complementary and velocity blocks exclude further
nonzero spectrum in a sufficiently small fixed neighborhood.
\end{proof}

\begin{remark}[Center, frequency, and linear time scale]
If $Q(-x)=Q(x)$, then $c=0$ by parity. Without this symmetry, omitting
$c^TT^{-1}c$ gives an incorrect coefficient in general.
Critical scaling also gives the exact spectral covariance
\[
 \lambda(\omega,v)=
 \omega\lambda\bigl(1,v\omega^{-(m-1)/m}\bigr).
\]
Thus for $m>2$ and fixed $\omega>0$,
\[
 \lambda_+(\omega,v)=\omega^{1/m}
 \sqrt{\frac{(m-2)v^TTv}{\mathcal D_Q}}+o(|v|).
\]
The reciprocal growth rate is of order $|v|^{-1}$. This appendix alone concerns linear exponential growth; Section~\ref{sec:r12-nonlinear} proves a nonlinear exit time in the specified quartic class. When $m=2$, the displayed leading coefficient vanishes; a general
homogeneous degree-two symbol need not be a quadratic form, so an exact
Galilean degeneracy is not asserted.
\end{remark}

\begin{corollary}[Critical Hartree interactions]
Theorem~\ref{thm:r11-spectral-pair} holds, with the same coefficient,
for the critical Hartree flow
\[
 iu_t=a(D)u-(K*|u|^2)u,
 \qquad K(x)=|x|^{-m},\quad 1<m<N,
\]
whenever its real ground wave and local branch satisfy the same four
hypotheses with
\[
 L_-=a(D)+1-(K*Q^2),\qquad
 L_+h=L_-h-2Q\bigl(K*(Qh)\bigr).
\]
It also holds for a positive-type anisotropic interaction of the same
homogeneity under those operator and regularity assumptions.
\end{corollary}
\begin{proof}
The proof above uses the nonlinearity through phase and translation
invariance, conjugation symmetry, critical scaling, and the stated
real-block operators. These properties hold for the Hartree energy.
In particular $L_+\Lambda Q=-mQ$, the parameter equations and the
charge-scaling identities are unchanged. No step in the resolvent
reduction uses locality of $L_+$. Thus the proof applies verbatim.
\end{proof}

\begin{remark}[Relation to established instability theory]
The general mechanism of a spectral pair leaving zero at charge
Hessian degeneracy belongs to established Hamiltonian bifurcation
theory. Here the additional calculation is its coefficient for
mass-critical translation branches, including the coupling to the
translation modes. The theorem proves spectral instability for
$m>2$ under its explicit assumptions. Nonlinear orbital instability
requires a sufficiently differentiable local flow and the associated
linear evolution bounds or unstable-manifold construction in a
specified phase space; it is not inferred here for every admissible
symbol and dimension. Likewise, real-profile nondegeneracy is not
proved by this section.
\end{remark}


\clearpage
\begin{thebibliography}{9}
\raggedright
\bibitem{BL} T. Boulenger and E. Lenzmann,
\emph{Blowup for biharmonic NLS}, Ann. Sci. ENS 50 (2017).
\href{https://arxiv.org/abs/1503.01741}{arXiv:1503.01741}.
\bibitem{BCGJ} D. Bonheure, J.-B. Cast\'eras, T. Gou and L. Jeanjean,
\emph{Strong instability of ground states to a fourth order Schr\"odinger equation},
IMRN 2019, 5299--5315.
\href{https://arxiv.org/abs/1703.07977}{arXiv:1703.07977}.
\bibitem{DM} L. Dinh, with Muse Spark,
\emph{Finite-Time Blow-Up of Radial Negative-Energy Solutions for the
Mass-Critical Biharmonic Nonlinear Schr\"odinger Equation}, manuscript,
2 October 2026.
\href{https://ai.meta.com/research/publications/finite-time-blow-up-of-radial-negative-energy-solutions-for-the-mass-critical-biharmonic-nonlinear-schrodinger-equation/}{Official manuscript and PDF}.
\bibitem{LS} E. Lenzmann and J. Sok,
\emph{A sharp rearrangement principle in Fourier space and symmetry
results for PDEs with arbitrary order}, IMRN 2021, 15040--15081.
\href{https://arxiv.org/abs/1805.06294}{arXiv:1805.06294}.
\bibitem{Gou} T. Gou,
\emph{Blowup of cylindrically symmetric solutions for biharmonic NLS},
\href{https://arxiv.org/abs/2205.08167}{arXiv:2205.08167}, version 3.
\bibitem{BLSS} L. Bugiera, E. Lenzmann, A. Schikorra and J. Sok,
\emph{On symmetry of traveling solitary waves for dispersion generalized NLS},
\href{https://arxiv.org/abs/1908.06771}{arXiv:1908.06771}.
\bibitem{LZ} Z. Lin and C. Zeng,
\emph{Instability, index theorem, and exponential trichotomy for linear
Hamiltonian PDEs},
\href{https://arxiv.org/abs/1703.04016}{arXiv:1703.04016}.
\bibitem{HKY} Y. Hong, S. Kwon and H. Yoon,
\emph{Global existence versus finite time blowup dichotomy for the system
of nonlinear Schr\"odinger equations},
\href{https://arxiv.org/abs/1712.06811}{arXiv:1712.06811}.
\end{thebibliography}
\end{document}
